\part{V--VI klasė}

\section{Natūralieji skaičiai, dalikliai ir kartotiniai}\label{sec:naturalieji-skaiciai-dalikliai-ir-kartotiniai}

\subsection{Apibrėžtys}
\begin{apibreztis}
Natūralieji skaičiai yra skaičiai, kuriais skaičiuojame objektus: \(1,2,3,\ldots\). Mokykloje dažnai nagrinėjame ir skaičių \(0\), nes jis reikalingas pozicinei dešimtainei sistemai, skaičių tiesei ir veiksmams aprašyti.
\end{apibreztis}
\begin{apibreztis}
Skaičius \(a\) yra skaičiaus \(b\) daliklis, jei \(b\) dalijasi iš \(a\) be liekanos. Tada rašome \(a\mid b\). Pavyzdžiui, \(6\mid 42\), nes \(42:6=7\).
\end{apibreztis}
\begin{apibreztis}
Skaičiaus \(a\) kartotiniai yra skaičiai, gaunami \(a\) dauginant iš natūraliųjų skaičių: \(a,2a,3a,\ldots\). Pavyzdžiui, skaičiaus \(8\) kartotiniai yra \(8,16,24,32,\ldots\).
\end{apibreztis}
\begin{apibreztis}
Pirminis skaičius yra natūralusis skaičius, didesnis už \(1\), turintis tik du daliklius: \(1\) ir save patį. Sudėtinis skaičius yra didesnis už \(1\) natūralusis skaičius, turintis daugiau kaip du daliklius.
\end{apibreztis}

\subsection{Teoremos ir savybės}
\begin{savybe}
Kiekvieną sudėtinį natūralųjį skaičių galima išskaidyti pirminių skaičių sandauga. Pavyzdžiui, \(60=2^2\cdot 3\cdot 5\).
\end{savybe}
\begin{savybe}
Jei skaičius dalijasi iš \(2\), jo paskutinis skaitmuo yra \(0,2,4,6\) arba \(8\). Jei skaičius dalijasi iš \(5\), jo paskutinis skaitmuo yra \(0\) arba \(5\). Jei skaičius dalijasi iš \(10\), jo paskutinis skaitmuo yra \(0\).
\end{savybe}
\begin{savybe}
Skaičius dalijasi iš \(3\), jei jo skaitmenų suma dalijasi iš \(3\); skaičius dalijasi iš \(9\), jei jo skaitmenų suma dalijasi iš \(9\).
\end{savybe}
\begin{savybe}
Jeigu \(a\mid b\) ir \(a\mid c\), tai \(a\mid (b+c)\) ir \(a\mid (b-c)\), kai \(b\ge c\). Todėl dalumo požymius galima taikyti ir tikrinant reiškinių reikšmes.
\end{savybe}

\subsection{Taisyklės ir algoritmai}
\textbf{Kaip rasti visus daliklius.}
\begin{enumerate}[label=\arabic*.]
\item Tikrinkite skaičius nuo \(1\) didėjančia tvarka.
\item Kiekvieną rastą daliklį rašykite poroje su kitu dalikliu, nes jei \(a\cdot b=n\), tai abu skaičiai yra \(n\) dalikliai.
\item Sustokite, kai poros pradeda kartotis; pakanka tikrinti iki skaičiaus, kurio kvadratas neviršija \(n\).
\item Surikiuokite daliklius didėjimo tvarka.
\end{enumerate}

\textbf{Kaip rasti DBD ir MBK skaidant pirminiais dauginamaisiais.}
\begin{enumerate}[label=\arabic*.]
\item Išskaidykite kiekvieną skaičių pirminiais dauginamaisiais.
\item DBD sudarykite iš bendrų pirminių dauginamųjų su mažiausiais laipsniais.
\item MBK sudarykite iš visų pasirodžiusių pirminių dauginamųjų su didžiausiais laipsniais.
\item Patikrinkite: \(\operatorname{DBD}(a,b)\cdot \operatorname{MBK}(a,b)=a\cdot b\), kai skaičiai teigiami.
\end{enumerate}

\paragraph{Dažnos klaidos.}
Skaičius \(1\) nėra pirminis, nes turi tik vieną daliklį. Dalikliai ir kartotiniai painiojami: \(6\) yra \(42\) daliklis, bet \(42\) yra \(6\) kartotinis. Randant MBK nepakanka sudauginti skaičių, jei jie turi bendrų daliklių: \(8\cdot 12=96\), bet \(\operatorname{MBK}(8,12)=24\).

\subsection{Iliustruojantys pavyzdžiai}
\textbf{1 pavyzdys.} Raskime visus skaičiaus \(36\) daliklius.

Rašome daliklių poras:
\[
1\cdot 36=36,\quad 2\cdot 18=36,\quad 3\cdot 12=36,\quad 4\cdot 9=36,\quad 6\cdot 6=36.
\]
Toliau tikrinti nereikia, nes po \(6\) poros kartotųsi. Visi dalikliai:
\[
1,2,3,4,6,9,12,18,36.
\]
\textbf{Atsakymas:} \(1,2,3,4,6,9,12,18,36\).

\textbf{2 pavyzdys.} Nustatykime, ar \(738\) dalijasi iš \(2\), \(3\), \(5\), \(9\).

Paskutinis skaitmuo yra \(8\), todėl skaičius dalijasi iš \(2\), bet nesidalija iš \(5\). Skaitmenų suma:
\[
7+3+8=18.
\]
Kadangi \(18\) dalijasi iš \(3\) ir iš \(9\), tai \(738\) dalijasi ir iš \(3\), ir iš \(9\).

\textbf{Atsakymas:} dalijasi iš \(2,3,9\), nesidalija iš \(5\).

\textbf{3 pavyzdys.} Raskime \(\operatorname{DBD}(84,126)\) ir \(\operatorname{MBK}(84,126)\).

Skaidome:
\[
84=2^2\cdot 3\cdot 7,\qquad 126=2\cdot 3^2\cdot 7.
\]
Bendri mažiausi laipsniai: \(2^1\), \(3^1\), \(7^1\), todėl
\[
\operatorname{DBD}(84,126)=2\cdot 3\cdot 7=42.
\]
Visi didžiausi laipsniai: \(2^2\), \(3^2\), \(7^1\), todėl
\[
\operatorname{MBK}(84,126)=2^2\cdot 3^2\cdot 7=252.
\]
\textbf{Atsakymas:} \(\operatorname{DBD}=42\), \(\operatorname{MBK}=252\).

\paragraph{Pasitikrinkite.}
\begin{enumerate}[label=\arabic*.]
\item Ar \(1245\) dalijasi iš \(3\), \(5\), \(9\)?
\item Išskaidykite \(90\) pirminiais dauginamaisiais.
\item Raskite \(\operatorname{DBD}(18,30)\) ir \(\operatorname{MBK}(18,30)\).
\end{enumerate}

\subsection{Ryšys su kitomis temomis}
Dalikliai ir kartotiniai padeda trumpinti trupmenas, rasti bendrą vardiklį, tirti skaičių sekas ir spręsti praktinius pasikartojimo uždavinius. Ši tema remiasi veiksmų su natūraliaisiais skaičiais savybėmis iš \ref{sec:aritmetiniai-veiksmai-ir-savybes} ir tiesiogiai reikalinga skyriams \ref{sec:paprastosios-ir-desimtaines-trupmenos} bei \ref{sec:dalumo-pozymiai-pirminiai-skaiciai-bdd-mbk}.

\section{Paprastosios ir dešimtainės trupmenos}\label{sec:paprastosios-ir-desimtaines-trupmenos}

\subsection{Apibrėžtys}
\begin{apibreztis}
Paprastoji trupmena \(\frac{a}{b}\) rodo, kad vienetas padalytas į \(b\) lygių dalių ir paimta \(a\) tokių dalių. Čia \(b\ne 0\).
\end{apibreztis}
\begin{apibreztis}
Skaitiklis yra trupmenos viršuje esantis skaičius, o vardiklis yra apačioje esantis skaičius. Vardiklis pasako, į kiek dalių padalytas vienetas, skaitiklis -- kiek dalių paimta.
\end{apibreztis}
\begin{apibreztis}
Taisyklingoji trupmena yra mažesnė už \(1\), todėl jos skaitiklis mažesnis už vardiklį. Netaisyklingoji trupmena yra lygi \(1\) arba didesnė už \(1\), todėl jos skaitiklis ne mažesnis už vardiklį.
\end{apibreztis}
\begin{apibreztis}
Dešimtainė trupmena yra skaičius su kableliu, kurio dešimtainės skiltys reiškia dešimtąsias, šimtąsias, tūkstantąsias ir taip toliau.
\end{apibreztis}

\subsection{Teoremos ir savybės}
\begin{savybe}
Trupmenos reikšmė nepasikeičia, jei jos skaitiklį ir vardiklį padauginame arba padalijame iš to paties nelygaus nuliui skaičiaus:
\[
\frac{a}{b}=\frac{a\cdot k}{b\cdot k},\qquad \frac{a}{b}=\frac{a:k}{b:k}.
\]
\end{savybe}
\begin{savybe}
Vienodų vardiklių trupmenos sudedamos ir atimamos sudedant arba atimant skaitiklius:
\[
\frac{a}{c}+\frac{b}{c}=\frac{a+b}{c},\qquad \frac{a}{c}-\frac{b}{c}=\frac{a-b}{c}.
\]
\end{savybe}
\begin{savybe}
Dešimtaines trupmenas sudedame ir atimame rašydami kablelį po kableliu, nes vienodos skiltys turi būti vienoje kolonoje.
\end{savybe}
\begin{savybe}
Trupmenos daugyba iš natūraliojo skaičiaus reiškia tokių pačių trupmenų sudėjimą:
\[
n\cdot \frac{a}{b}=\frac{n\cdot a}{b}.
\]
\end{savybe}

\subsection{Taisyklės ir algoritmai}
\textbf{Trupmenos trumpinimas.}
\begin{enumerate}[label=\arabic*.]
\item Raskite bendrą skaitiklio ir vardiklio daliklį.
\item Padalykite ir skaitiklį, ir vardiklį iš to paties daliklio.
\item Kartokite, kol skaitiklis ir vardiklis nebeturi bendrų daliklių, didesnių už \(1\).
\end{enumerate}

\textbf{Skirtingų vardiklių trupmenų sudėtis arba atimtis.}
\begin{enumerate}[label=\arabic*.]
\item Raskite bendrą vardiklį, dažnai patogiausia -- MBK.
\item Išplėskite kiekvieną trupmeną iki bendro vardiklio.
\item Sudėkite arba atimkite skaitiklius.
\item Jei įmanoma, sutrumpinkite rezultatą arba išskirkite sveikąją dalį.
\end{enumerate}

\paragraph{Dažnos klaidos.}
Negalima sudėti vardiklių: \(\frac{1}{3}+\frac{1}{4}\ne \frac{2}{7}\). Trumpinant reikia dalyti ir skaitiklį, ir vardiklį. Lyginant dešimtainius skaičius verta prirašyti nulių: \(2,7=2,70\), todėl \(2,7>2,68\).

\subsection{Iliustruojantys pavyzdžiai}
\textbf{1 pavyzdys.} Sutrumpinkime \(\frac{18}{24}\).

Skaičiai \(18\) ir \(24\) dalijasi iš \(6\):
\[
\frac{18}{24}=\frac{18:6}{24:6}=\frac{3}{4}.
\]
\textbf{Atsakymas:} \(\frac{3}{4}\).

\textbf{2 pavyzdys.} Apskaičiuokime \(\frac{2}{5}+\frac{1}{10}\).

Bendras vardiklis yra \(10\). Pirmą trupmeną išplečiame:
\[
\frac{2}{5}=\frac{4}{10}.
\]
Tada
\[
\frac{4}{10}+\frac{1}{10}=\frac{5}{10}=\frac{1}{2}.
\]
\textbf{Atsakymas:} \(\frac{1}{2}\).

\textbf{3 pavyzdys.} Mišrųjį skaičių \(3\frac{2}{5}\) užrašykime dešimtaine trupmena.

Trupmeninę dalį keičiame:
\[
\frac{2}{5}=\frac{4}{10}=0,4.
\]
Todėl
\[
3\frac{2}{5}=3,4.
\]
\textbf{Atsakymas:} \(3,4\).

\textbf{4 pavyzdys.} Apskaičiuokime \(4,06+12,8-3,175\).

Sulyginame skiltis:
\[
4,060+12,800=16,860,\qquad 16,860-3,175=13,685.
\]
\textbf{Atsakymas:} \(13,685\).

\paragraph{Pasitikrinkite.}
\begin{enumerate}[label=\arabic*.]
\item Sutrumpinkite \(\frac{28}{42}\).
\item Apskaičiuokite \(\frac{3}{4}-\frac{1}{6}\).
\item Kuris skaičius didesnis: \(5,09\) ar \(5,1\)?
\end{enumerate}

\subsection{Ryšys su kitomis temomis}
Trupmenos remiasi dalikliais ir kartotiniais iš \ref{sec:naturalieji-skaiciai-dalikliai-ir-kartotiniai}. Jos būtinos procentams \ref{sec:procentai-santykiai-ir-proporcijos-pradzia}, racionaliesiems skaičiams \ref{sec:sveikieji-skaiciai-ir-racionalieji-skaiciai}, tikimybei \ref{sec:atsitiktiniai-bandymas-ir-tikimybe} ir masteliui \ref{sec:plokstumos-ir-erdves-figuru-konstravimas}.

\section{Procentai, santykiai ir proporcijos pradžia}\label{sec:procentai-santykiai-ir-proporcijos-pradzia}

\subsection{Apibrėžtys}
\begin{apibreztis}
Procentas yra viena šimtoji dalis: \(1\%=\frac{1}{100}=0,01\). Todėl \(37\%\) reiškia \(37\) šimtąsias.
\end{apibreztis}
\begin{apibreztis}
Santykis palygina du dydžius dalybos būdu. Santykis \(3:5\) reiškia, kad pirmajam dydžiui tenka \(3\) dalys, o antrajam -- \(5\) dalys.
\end{apibreztis}
\begin{apibreztis}
Proporcija yra dviejų santykių lygybė, pavyzdžiui, \(\frac{a}{b}=\frac{c}{d}\), kai \(b\ne 0\) ir \(d\ne 0\).
\end{apibreztis}
\begin{apibreztis}
Procentinis pokytis rodo, keliais procentais dydis padidėjo arba sumažėjo, lyginant su pradine reikšme.
\end{apibreztis}

\subsection{Teoremos ir savybės}
\begin{savybe}
Norėdami rasti \(p\%\) skaičiaus \(a\), skaičiuojame \(\frac{p}{100}\cdot a\).
\end{savybe}
\begin{savybe}
Jei dalis sudaro \(p\%\) visumos, tai visumą galima rasti dalį dalijant iš \(\frac{p}{100}\).
\end{savybe}
\begin{savybe}
Proporcijoje \(\frac{a}{b}=\frac{c}{d}\) galioja kryžminė sandauga:
\[
a\cdot d=b\cdot c.
\]
\end{savybe}
\begin{savybe}
Santykį galima trumpinti arba plėsti taip pat kaip trupmeną, jei abu santykio nariai dauginami arba dalijami iš to paties nelygaus nuliui skaičiaus.
\end{savybe}

\subsection{Taisyklės ir algoritmai}
\textbf{Procentų keitimas.}
\begin{enumerate}[label=\arabic*.]
\item Procentus į dešimtainę trupmeną keiskite dalydami iš \(100\): \(35\%=0,35\).
\item Dešimtainę trupmeną į procentus keiskite daugindami iš \(100\): \(0,07=7\%\).
\item Paprastąsias trupmenas patogu paversti procentais radus lygiavertę trupmeną su vardikliu \(100\), jei tai įmanoma.
\end{enumerate}

\textbf{Proporcijos uždavinys.}
\begin{enumerate}[label=\arabic*.]
\item Aiškiai įvardykite, kurie dydžiai lyginami.
\item Sudarykite santykių lygybę.
\item Taikykite kryžminę sandaugą arba raskite vieneto reikšmę.
\item Patikrinkite, ar atsakymas turi tinkamus vienetus ir realią prasmę.
\end{enumerate}

\paragraph{Dažnos klaidos.}
Procentai visada skaičiuojami nuo konkrečios visumos. \(20\%\) nuo \(50\) ir \(20\%\) nuo \(80\) nėra tas pats skaičius. Nuolaida mažina pradinę kainą, todėl po \(30\%\) nuolaidos lieka \(70\%\) kainos, o ne \(30\%\).

\subsection{Iliustruojantys pavyzdžiai}
\textbf{1 pavyzdys.} Raskime \(20\%\) skaičiaus \(45\).
\[
20\%=0,20,\qquad 0,20\cdot 45=9.
\]
\textbf{Atsakymas:} \(9\).

\textbf{2 pavyzdys.} Supaprastinkime santykį \(18:24\).
\[
18:24=(18:6):(24:6)=3:4.
\]
\textbf{Atsakymas:} \(3:4\).

\textbf{3 pavyzdys.} Jei \(4\) sąsiuviniai kainuoja \(2,40\) Eur, kiek kainuoja \(7\) tokie sąsiuviniai?

Vieno sąsiuvinio kaina:
\[
2,40:4=0,60\text{ Eur}.
\]
Septyni sąsiuviniai kainuoja:
\[
7\cdot 0,60=4,20\text{ Eur}.
\]
\textbf{Atsakymas:} \(4,20\) Eur.

\textbf{4 pavyzdys.} Prekė kainavo \(80\) Eur. Kaina sumažinta \(15\%\). Kokia nauja kaina?

Nuolaida:
\[
0,15\cdot 80=12\text{ Eur}.
\]
Nauja kaina:
\[
80-12=68\text{ Eur}.
\]
\textbf{Atsakymas:} \(68\) Eur.

\paragraph{Pasitikrinkite.}
\begin{enumerate}[label=\arabic*.]
\item Raskite \(12\%\) skaičiaus \(250\).
\item Santykį \(45:60\) užrašykite paprasčiausiu pavidalu.
\item \(6\) kg obuolių kainuoja \(9\) Eur. Kiek kainuoja \(10\) kg?
\end{enumerate}

\subsection{Ryšys su kitomis temomis}
Procentai yra trupmenų taikymas, todėl remiasi \ref{sec:paprastosios-ir-desimtaines-trupmenos}. Santykiai ir proporcijos padeda nagrinėti tiesioginį proporcingumą \ref{sec:tiesiniai-ir-netiesiniai-sarysiai-5-6}, mastelį, finansinius skaičiavimus \ref{sec:finansai-biudzetas-nuolaidos-ir-palukanos} ir tikimybes.

\section{Sveikieji skaičiai ir racionalieji skaičiai}\label{sec:sveikieji-skaiciai-ir-racionalieji-skaiciai}

\subsection{Apibrėžtys}
\begin{apibreztis}
Sveikieji skaičiai yra \(\ldots,-3,-2,-1,0,1,2,3,\ldots\). Teigiamieji sveikieji skaičiai yra dešinėje nuo nulio, neigiamieji -- kairėje.
\end{apibreztis}
\begin{apibreztis}
Priešingi skaičiai yra vienodai nutolę nuo nulio, bet yra skirtingose jo pusėse, pavyzdžiui, \(8\) ir \(-8\). Jų suma lygi \(0\).
\end{apibreztis}
\begin{apibreztis}
Skaičiaus modulis \(|a|\) yra atstumas nuo skaičiaus \(a\) iki nulio skaičių tiesėje. Todėl \(|-6|=6\) ir \(|6|=6\).
\end{apibreztis}
\begin{apibreztis}
Racionalusis skaičius yra skaičius, kurį galima užrašyti trupmena \(\frac{a}{b}\), kur \(a\) ir \(b\) yra sveikieji skaičiai, o \(b\ne 0\).
\end{apibreztis}

\subsection{Teoremos ir savybės}
\begin{savybe}
Skaičius, esantis skaičių tiesėje dešiniau, yra didesnis. Todėl \(-2>-5\), nes \(-2\) yra arčiau nulio ir dešiniau.
\end{savybe}
\begin{savybe}
Priešingų skaičių suma lygi nuliui: \(a+(-a)=0\).
\end{savybe}
\begin{savybe}
Sudauginus arba padalijus vienodo ženklo skaičius gaunamas teigiamas skaičius; sudauginus arba padalijus skirtingų ženklų skaičius gaunamas neigiamas skaičius.
\end{savybe}
\begin{savybe}
Kiekviena paprastoji trupmena gali būti užrašyta baigtine arba begaline periodine dešimtaine trupmena.
\end{savybe}

\subsection{Taisyklės ir algoritmai}
\textbf{Sveikųjų skaičių sudėtis.}
\begin{enumerate}[label=\arabic*.]
\item Jei ženklai vienodi, sudėkite modulius ir palikite bendrą ženklą.
\item Jei ženklai skirtingi, iš didesnio modulio atimkite mažesnį.
\item Rezultatui palikite didesnio modulio skaičiaus ženklą.
\end{enumerate}

\textbf{Racionaliųjų skaičių veiksmai.}
\begin{enumerate}[label=\arabic*.]
\item Pirmiausia sutvarkykite ženklus.
\item Trupmenoms taikykite žinomas trupmenų taisykles: bendras vardiklis sudėčiai, skaitiklių ir vardiklių daugyba daugybai, atvirkštinis skaičius dalybai.
\item Atsakymą supaprastinkite ir, jei reikia, suapvalinkite nurodytu tikslumu.
\end{enumerate}

\paragraph{Dažnos klaidos.}
Neigiamas skaičius su didesniu moduliu yra mažesnis: \(-10<-3\). Ženklas „minus“ prieš trupmeną gali būti rašomas prie skaitiklio, prie vardiklio arba prieš visą trupmeną, bet ne dviejose vietose iš karto: \(-\frac{2}{5}=\frac{-2}{5}=\frac{2}{-5}\), o \(\frac{-2}{-5}=\frac{2}{5}\).

\subsection{Iliustruojantys pavyzdžiai}
\textbf{1 pavyzdys.} Apskaičiuokime \(-7+12\).

Ženklai skirtingi, todėl atimame modulius:
\[
12-7=5.
\]
Didesnį modulį turi \(12\), todėl rezultatas teigiamas.

\textbf{Atsakymas:} \(5\).

\textbf{2 pavyzdys.} Apskaičiuokime \(-4\cdot 6\).
\[
4\cdot 6=24.
\]
Ženklai skirtingi, todėl sandauga neigiama.

\textbf{Atsakymas:} \(-24\).

\textbf{3 pavyzdys.} Palyginkime \(-\frac{2}{3}\) ir \(-\frac{1}{2}\).

Teigiamos trupmenos:
\[
\frac{2}{3}=\frac{4}{6},\qquad \frac{1}{2}=\frac{3}{6},
\]
todėl \(\frac{2}{3}>\frac{1}{2}\). Neigiamųjų skaičių tvarka priešinga:
\[
-\frac{2}{3}<-\frac{1}{2}.
\]
\textbf{Atsakymas:} \(-\frac{2}{3}<-\frac{1}{2}\).

\textbf{4 pavyzdys.} Apskaičiuokime \(-\frac{3}{4}+\frac{5}{6}\).

Bendras vardiklis yra \(12\):
\[
-\frac{3}{4}=-\frac{9}{12},\qquad \frac{5}{6}=\frac{10}{12}.
\]
Tada
\[
-\frac{9}{12}+\frac{10}{12}=\frac{1}{12}.
\]
\textbf{Atsakymas:} \(\frac{1}{12}\).

\paragraph{Pasitikrinkite.}
\begin{enumerate}[label=\arabic*.]
\item Apskaičiuokite \(-15+9\).
\item Apskaičiuokite \((-8):(-2)\).
\item Palyginkite \(-0,4\) ir \(-\frac{1}{3}\).
\end{enumerate}

\subsection{Ryšys su kitomis temomis}
Sveikieji ir racionalieji skaičiai išplečia natūraliųjų skaičių ir trupmenų pasaulį. Jie reikalingi sprendžiant lygtis \ref{sec:lygtys-ir-nelygybes-su-vienu-nezinomuoju}, žymint taškus koordinačių plokštumoje \ref{sec:dydziu-priklausomybes-ir-koordinaciu-plokstuma} ir vėliau mokantis realiųjų skaičių.

\section{Skaitiniai ir raidiniai reiškiniai}\label{sec:skaitiniai-ir-raidiniai-reiskiniai}

\subsection{Apibrėžtys}
\begin{apibreztis}
Skaitinis reiškinys yra skaičių, veiksmų ženklų ir skliaustų užrašas, turintis reikšmę, pavyzdžiui, \(18-2\cdot(3+4)\).
\end{apibreztis}
\begin{apibreztis}
Raidinis reiškinys yra reiškinys, kuriame yra raidžių, žyminčių skaičius, pavyzdžiui, \(2x+5\).
\end{apibreztis}
\begin{apibreztis}
Reiškinio reikšmė gaunama vietoje raidžių įrašius skaičius ir atlikus veiksmus.
\end{apibreztis}
\begin{apibreztis}
Panašieji nariai yra raidinio reiškinio nariai, turintys tą pačią raidinę dalį, pavyzdžiui, \(3x\) ir \(-5x\).
\end{apibreztis}

\subsection{Teoremos ir savybės}
\begin{savybe}
Veiksmų tvarka: pirmiausia atliekami veiksmai skliaustuose, tada daugyba ir dalyba, galiausiai sudėtis ir atimtis iš kairės į dešinę.
\end{savybe}
\begin{savybe}
Panašiuosius narius galima sudėti sudedant jų koeficientus:
\[
3x+2x=5x,\qquad 7a-4a=3a.
\]
\end{savybe}
\begin{savybe}
Skliaustų atskleidimas remiasi skirstomumo savybe:
\[
a(b+c)=ab+ac,\qquad a(b-c)=ab-ac.
\]
\end{savybe}
\begin{savybe}
Jei prieš skliaustus yra minusas, nuimant skliaustus visų narių ženklai skliaustuose pasikeičia priešingais.
\end{savybe}

\subsection{Taisyklės ir algoritmai}
\textbf{Reiškinio reikšmės radimas.}
\begin{enumerate}[label=\arabic*.]
\item Vietoje raidės įrašykite duotą skaičių į visas tos raidės vietas.
\item Jei raidė dauginama iš skaičiaus, užrašykite daugybos ženklą, kad nesupainiotumėte, pavyzdžiui, \(3x=3\cdot x\).
\item Atlikite veiksmus pagal veiksmų tvarką.
\item Jei skaičiai neigiami arba trupmeniniai, ypač atidžiai tikrinkite ženklus.
\end{enumerate}

\textbf{Reiškinio prastinimas.}
\begin{enumerate}[label=\arabic*.]
\item Atskleiskite skliaustus, jei tai patogu.
\item Suraskite panašiuosius narius.
\item Sudėkite jų koeficientus.
\item Užrašykite paprasčiausią reiškinio pavidalą.
\end{enumerate}

\paragraph{Dažnos klaidos.}
Užrašas \(2x\) reiškia \(2\cdot x\), o ne skaičių, kurio dešimčių skaitmuo \(2\). Skliaustai po minuso keičia visus ženklus: \(-(x-5)=-x+5\). Panašūs yra tik tie nariai, kurių raidinė dalis sutampa: \(3x\) ir \(3y\) nesutraukiami.

\subsection{Iliustruojantys pavyzdžiai}
\textbf{1 pavyzdys.} Apskaičiuokime \(18-2\cdot(3+4)\).
\[
3+4=7,\qquad 2\cdot 7=14,\qquad 18-14=4.
\]
\textbf{Atsakymas:} \(4\).

\textbf{2 pavyzdys.} Raskime \(3x+7\) reikšmę, kai \(x=5\).
\[
3x+7=3\cdot 5+7=15+7=22.
\]
\textbf{Atsakymas:} \(22\).

\textbf{3 pavyzdys.} Supaprastinkime \(4a+3+2a-1\).
\[
4a+2a=6a,\qquad 3-1=2,
\]
todėl
\[
4a+3+2a-1=6a+2.
\]
\textbf{Atsakymas:} \(6a+2\).

\textbf{4 pavyzdys.} Atskleiskime skliaustus ir sutraukime narius: \(3(2x-5)-4x\).
\[
3(2x-5)-4x=6x-15-4x=2x-15.
\]
\textbf{Atsakymas:} \(2x-15\).

\paragraph{Pasitikrinkite.}
\begin{enumerate}[label=\arabic*.]
\item Apskaičiuokite \(24:3+2\cdot 5\).
\item Raskite \(5a-2\) reikšmę, kai \(a=4\).
\item Supaprastinkite \(7x-3+2x+8\).
\end{enumerate}

\subsection{Ryšys su kitomis temomis}
Reiškiniai yra algebrai reikalinga kalba. Jie remiasi veiksmų savybėmis iš \ref{sec:aritmetiniai-veiksmai-ir-savybes}, padeda spręsti lygtis \ref{sec:lygtys-ir-nelygybes-su-vienu-nezinomuoju} ir aprašyti dėsningumus \ref{sec:tiesiniai-ir-netiesiniai-sarysiai-5-6}.

\section{Lygtys ir nelygybės su vienu nežinomuoju}\label{sec:lygtys-ir-nelygybes-su-vienu-nezinomuoju}

\subsection{Apibrėžtys}
\begin{apibreztis}
Lygtis yra lygybė su nežinomuoju, pavyzdžiui, \(x+5=12\).
\end{apibreztis}
\begin{apibreztis}
Lygties sprendinys yra nežinomojo reikšmė, su kuria lygybė tampa teisinga.
\end{apibreztis}
\begin{apibreztis}
Nelygybė yra užrašas su ženklais \(<\), \(>\), \(\le\) arba \(\ge\), pavyzdžiui, \(x+2<7\).
\end{apibreztis}
\begin{apibreztis}
Sprendinių aibė yra visi skaičiai, kurie tinka lygčiai arba nelygybei.
\end{apibreztis}

\subsection{Teoremos ir savybės}
\begin{savybe}
Prie abiejų lygties pusių galima pridėti tą patį skaičių arba iš abiejų pusių atimti tą patį skaičių; sprendiniai nepasikeičia.
\end{savybe}
\begin{savybe}
Abi lygties puses galima padauginti arba padalyti iš to paties nelygaus nuliui skaičiaus; sprendiniai nepasikeičia.
\end{savybe}
\begin{savybe}
Sprendžiant nelygybę, pridėjus ar atėmus tą patį skaičių iš abiejų pusių, nelygybės ženklas nesikeičia.
\end{savybe}
\begin{savybe}
Dauginant arba dalijant nelygybę iš neigiamo skaičiaus, nelygybės ženklas apsiverčia.
\end{savybe}

\subsection{Taisyklės ir algoritmai}
\textbf{Lygties sprendimas.}
\begin{enumerate}[label=\arabic*.]
\item Jei yra skliaustų, atskleiskite juos.
\item Nežinomojo narius perkelkite į vieną pusę, skaičius -- į kitą.
\item Atlikite veiksmus su panašiaisiais nariais.
\item Padalykite iš nežinomojo koeficiento.
\item Patikrinkite sprendinį įrašydami jį į pradinę lygtį.
\end{enumerate}

\textbf{Tekstinio uždavinio lygtis.}
\begin{enumerate}[label=\arabic*.]
\item Pasirinkite, ką žymės nežinomasis.
\item Iš sąlygos sudarykite dvi to paties dydžio išraiškas.
\item Parašykite lygtį ir ją išspręskite.
\item Atsakymą užrašykite sakiniais su vienetais.
\end{enumerate}

\paragraph{Dažnos klaidos.}
Lygties pusiausvyra išlaikoma tik tada, kai tas pats veiksmas atliekamas abiejose pusėse. Nelygybės ženklas apsiverčia tik dauginant arba dalijant iš neigiamo skaičiaus, ne tada, kai pridedame neigiamą skaičių. Patikrinimas turi būti atliekamas pradinėje, o ne jau pertvarkytoje lygtyje.

\subsection{Iliustruojantys pavyzdžiai}
\textbf{1 pavyzdys.} Išspręskime lygtį \(x+8=21\).
\[
x+8-8=21-8,\qquad x=13.
\]
Patikrinimas: \(13+8=21\).

\textbf{Atsakymas:} \(x=13\).

\textbf{2 pavyzdys.} Išspręskime lygtį \(3x-4=17\).
\[
3x=21,\qquad x=7.
\]
Patikrinimas: \(3\cdot 7-4=17\).

\textbf{Atsakymas:} \(x=7\).

\textbf{3 pavyzdys.} Išspręskime nelygybę \(2x+1<9\).
\[
2x<8,\qquad x<4.
\]
\textbf{Atsakymas:} visi skaičiai, mažesni už \(4\).

\textbf{4 pavyzdys.} Į klasės ekskursiją važiuoja \(24\) mokiniai. Bilietas mokiniui kainuoja \(x\) Eur, o bendra bilietų kaina \(144\) Eur. Raskime bilieto kainą.
\[
24x=144,\qquad x=144:24=6.
\]
\textbf{Atsakymas:} vienas bilietas kainuoja \(6\) Eur.

\paragraph{Pasitikrinkite.}
\begin{enumerate}[label=\arabic*.]
\item Išspręskite \(x-9=14\).
\item Išspręskite \(5x+2=37\).
\item Išspręskite \(3x-6\ge 12\).
\end{enumerate}

\subsection{Ryšys su kitomis temomis}
Lygtys ir nelygybės naudoja raidinius reiškinius iš \ref{sec:skaitiniai-ir-raidiniai-reiskiniai}. Jos padeda aprašyti dydžių priklausomybes \ref{sec:dydziu-priklausomybes-ir-koordinaciu-plokstuma}, procentinius uždavinius ir geometrijos situacijas.

\section{Dydžių priklausomybės ir koordinačių plokštuma}\label{sec:dydziu-priklausomybes-ir-koordinaciu-plokstuma}

\subsection{Apibrėžtys}
\begin{apibreztis}
Dydžių priklausomybė rodo, kaip vieno dydžio reikšmė keičiasi, kai keičiasi kitas dydis.
\end{apibreztis}
\begin{apibreztis}
Koordinačių plokštumą sudaro dvi statmenos ašys: horizontalioji \(x\) ašis ir vertikalioji \(y\) ašis. Jų susikirtimo taškas vadinamas pradžios tašku.
\end{apibreztis}
\begin{apibreztis}
Taško koordinatės užrašomos pora \((x;y)\): pirmasis skaičius rodo poslinkį horizontaliai, antrasis -- vertikaliai.
\end{apibreztis}
\begin{apibreztis}
Grafikas yra priklausomybės vaizdas koordinačių plokštumoje. Jis padeda pamatyti, kaip kinta dydžiai.
\end{apibreztis}

\subsection{Teoremos ir savybės}
\begin{savybe}
Jei dydžiai tiesiogiai proporcingi, padidinus vieną dydį kelis kartus, tiek pat kartų padidėja ir kitas. Tokią priklausomybę galima rašyti \(y=kx\), kur \(k\) yra pastovus santykis.
\end{savybe}
\begin{savybe}
Taškas ant \(x\) ašies turi koordinatę \(y=0\), o taškas ant \(y\) ašies turi koordinatę \(x=0\).
\end{savybe}
\begin{savybe}
Lentelė, žodinis aprašas, formulė ir grafikas gali vaizduoti tą pačią priklausomybę.
\end{savybe}
\begin{savybe}
Koordinačių ašys nepriklauso nė vienam ketvirčiui. Ketvirčiai numeruojami prieš laikrodžio rodyklę, pradedant nuo srities, kur \(x>0\) ir \(y>0\).
\end{savybe}

\subsection{Taisyklės ir algoritmai}
\textbf{Taško žymėjimas.}
\begin{enumerate}[label=\arabic*.]
\item Pradėkite nuo pradžios taško \((0;0)\).
\item Pagal pirmą koordinatę judėkite į dešinę, jei ji teigiama, arba į kairę, jei neigiama.
\item Pagal antrą koordinatę judėkite aukštyn, jei ji teigiama, arba žemyn, jei neigiama.
\item Pažymėkite tašką ir užrašykite jo vardą.
\end{enumerate}

\textbf{Priklausomybės lentelė ir grafikas.}
\begin{enumerate}[label=\arabic*.]
\item Pasirinkite kelias \(x\) reikšmes.
\item Pagal taisyklę apskaičiuokite \(y\).
\item Taškus \((x;y)\) pažymėkite koordinačių plokštumoje.
\item Jei dydžiai kinta tolygiai, taškus galima jungti linija; jei kalbama apie atskirus objektus, taškai paliekami atskiri.
\end{enumerate}

\paragraph{Dažnos klaidos.}
Koordinatėse svarbi tvarka: \((2;5)\) ir \((5;2)\) yra skirtingi taškai. Ne kiekviena lentelė rodo tiesioginį proporcingumą: santykis \(y:x\) turi būti pastovus. Jei \(x=0\), santykio \(y:x\) skaičiuoti negalima.

\subsection{Iliustruojantys pavyzdžiai}
\textbf{1 pavyzdys.} Kiek kainuoja \(1\), \(2\), \(3\), \(4\) pieštukai, jei vienas kainuoja \(0,30\) Eur?
\[
y=0,30x.
\]
\[
\begin{array}{c|cccc}
x&1&2&3&4\\ \hline
y&0,30&0,60&0,90&1,20
\end{array}
\]
\textbf{Atsakymas:} \(0,30\), \(0,60\), \(0,90\), \(1,20\) Eur.

\textbf{2 pavyzdys.} Nustatykime, kuriame ketvirtyje yra taškas \(A(-3;2)\).

Pirmoji koordinatė neigiama, todėl taškas kairėje nuo \(y\) ašies. Antroji koordinatė teigiama, todėl taškas virš \(x\) ašies. Tai II ketvirtis.

\textbf{Atsakymas:} II ketvirtis.

\textbf{3 pavyzdys.} Pagal taisyklę \(y=2x+1\) raskime \(y\), kai \(x=0,2,5\).
\[
x=0:\ y=1,\qquad x=2:\ y=5,\qquad x=5:\ y=11.
\]
\textbf{Atsakymas:} gauname taškus \((0;1)\), \((2;5)\), \((5;11)\).

\paragraph{Pasitikrinkite.}
\begin{enumerate}[label=\arabic*.]
\item Kuriame ketvirtyje yra taškas \(B(4;-6)\)?
\item Užpildykite lentelę \(y=3x\), kai \(x=1,2,5\).
\item Ar lentelė \((1;4)\), \((2;8)\), \((3;12)\) rodo tiesioginį proporcingumą?
\end{enumerate}

\subsection{Ryšys su kitomis temomis}
Dydžių priklausomybės siejasi su proporcijomis iš \ref{sec:procentai-santykiai-ir-proporcijos-pradzia}, lygtimis iš \ref{sec:lygtys-ir-nelygybes-su-vienu-nezinomuoju} ir vėliau tampa tiesinių funkcijų pagrindu.

\section{Kampai, trikampiai ir keturkampiai}\label{sec:kampai-trikampiai-ir-keturkampiai}

\subsection{Apibrėžtys}
\begin{apibreztis}
Kampas yra figūra, sudaryta iš dviejų spindulių, turinčių bendrą pradžios tašką. Kampo dydis matuojamas laipsniais.
\end{apibreztis}
\begin{apibreztis}
Trikampis yra daugiakampis, turintis tris kraštines, tris viršūnes ir tris kampus.
\end{apibreztis}
\begin{apibreztis}
Keturkampis yra daugiakampis, turintis keturias kraštines, keturias viršūnes ir keturis kampus.
\end{apibreztis}
\begin{apibreztis}
Gretutiniai kampai turi vieną bendrą kraštinę, o kitos jų kraštinės sudaro tiesę. Kryžminiai kampai susidaro susikertant dviem tiesėms ir yra vienas prieš kitą.
\end{apibreztis}

\subsection{Teoremos ir savybės}
\begin{teorema}
Trikampio kampų suma lygi \(180^\circ\).
\end{teorema}
\begin{teorema}
Keturkampio kampų suma lygi \(360^\circ\).
\end{teorema}
\begin{savybe}
Stačiojo kampo dydis yra \(90^\circ\), ištiestinio kampo -- \(180^\circ\), pilnojo kampo -- \(360^\circ\).
\end{savybe}
\begin{savybe}
Gretutinių kampų suma lygi \(180^\circ\), o kryžminiai kampai yra lygūs.
\end{savybe}
\begin{savybe}
Stačiakampio priešingos kraštinės yra lygios, visi kampai statieji. Rombas turi visas kraštines lygias, o lygiagretainio priešingos kraštinės yra lygiagrečios ir lygios.
\end{savybe}

\subsection{Taisyklės ir algoritmai}
\textbf{Kampo matavimas matlankiu.}
\begin{enumerate}[label=\arabic*.]
\item Matlankio centrą dėkite ant kampo viršūnės.
\item Vieną kampo kraštinę sutapatinkite su \(0^\circ\) linija.
\item Skaitykite tą skalę, kuri prasideda nuo šios \(0^\circ\) linijos.
\item Įvertinkite, ar gautas kampas pagal vaizdą yra smailusis, statusis ar bukasis.
\end{enumerate}

\textbf{Nežinomo kampo radimas.}
\begin{enumerate}[label=\arabic*.]
\item Nuspręskite, kuri savybė tinka: trikampio kampų suma, keturkampio kampų suma, gretutiniai ar kryžminiai kampai.
\item Sudarykite lygtį arba atlikite atimtį.
\item Patikrinkite, ar kampo dydis atitinka figūros rūšį.
\end{enumerate}

\paragraph{Dažnos klaidos.}
Matlankyje yra dvi skalės, todėl reikia pasirinkti tą, kuri prasideda nuo kampo kraštinės. Trikampio kampų suma visada \(180^\circ\), net jei trikampis nubrėžtas didelis. Gretutiniai kampai ne būtinai lygūs, bet jų suma lygi \(180^\circ\).

\subsection{Iliustruojantys pavyzdžiai}
\textbf{1 pavyzdys.} Trikampio du kampai yra \(48^\circ\) ir \(67^\circ\). Raskime trečią kampą.
\[
180^\circ-48^\circ-67^\circ=65^\circ.
\]
\textbf{Atsakymas:} \(65^\circ\).

\textbf{2 pavyzdys.} Vienas iš gretutinių kampų yra \(112^\circ\). Raskime kitą.
\[
180^\circ-112^\circ=68^\circ.
\]
\textbf{Atsakymas:} \(68^\circ\).

\textbf{3 pavyzdys.} Keturkampio trys kampai yra \(90^\circ\), \(120^\circ\) ir \(75^\circ\). Raskime ketvirtą.
\[
360^\circ-(90^\circ+120^\circ+75^\circ)=75^\circ.
\]
\textbf{Atsakymas:} \(75^\circ\).

\paragraph{Pasitikrinkite.}
\begin{enumerate}[label=\arabic*.]
\item Ar gali trikampio kampai būti \(40^\circ\), \(60^\circ\), \(90^\circ\)?
\item Vienas kryžminis kampas yra \(35^\circ\). Koks jam kryžminis kampas?
\item Stačiakampio ilgis \(8\) cm, plotis \(5\) cm. Raskite perimetrą.
\end{enumerate}

\subsection{Ryšys su kitomis temomis}
Kampų savybės reikalingos braižymui \ref{sec:plokstumos-ir-erdves-figuru-konstravimas}, plotams \ref{sec:plotas-pavirsius-ir-turis} ir įrodymų supratimui. Ši tema remiasi ankstesniu figūrų pažinimu ir vėliau veda į trikampių lygumą bei panašumą.

\section{Apskritimas, skritulys ir pi}\label{sec:apskritimas-skritulys-ir-pi}

\subsection{Apibrėžtys}
\begin{apibreztis}
Apskritimas yra visų plokštumos taškų, vienodai nutolusių nuo vieno taško, vadinamo centru, aibė.
\end{apibreztis}
\begin{apibreztis}
Skritulys yra apskritimo apribota plokštumos dalis kartu su pačiu apskritimu.
\end{apibreztis}
\begin{apibreztis}
Spindulys \(r\) jungia centrą su apskritimo tašku. Skersmuo \(d\) jungia du apskritimo taškus per centrą, todėl \(d=2r\).
\end{apibreztis}
\begin{apibreztis}
Skaičius \(\pi\) yra apskritimo ilgio ir skersmens santykis. Apytiksliai \(\pi\approx 3,14\).
\end{apibreztis}

\subsection{Teoremos ir savybės}
\begin{savybe}
Apskritimo ilgis apskaičiuojamas formule
\[
C=2\pi r=\pi d.
\]
\end{savybe}
\begin{savybe}
Skritulio plotas apskaičiuojamas formule
\[
S=\pi r^2.
\]
\end{savybe}
\begin{savybe}
Visi to paties apskritimo spinduliai yra lygūs.
\end{savybe}
\begin{savybe}
Jei spindulys padidėja \(2\) kartus, apskritimo ilgis padidėja \(2\) kartus, o skritulio plotas -- \(4\) kartus.
\end{savybe}

\subsection{Taisyklės ir algoritmai}
\textbf{Uždaviniai su apskritimu.}
\begin{enumerate}[label=\arabic*.]
\item Perskaitykite, ar duotas spindulys, ar skersmuo.
\item Jei duotas skersmuo, raskite spindulį: \(r=\frac{d}{2}\).
\item Pasirinkite formulę: ilgiui \(C=2\pi r\), plotui \(S=\pi r^2\).
\item Jei prašoma apytikslio atsakymo, naudokite \(\pi\approx 3,14\) ir suapvalinkite.
\end{enumerate}

\paragraph{Dažnos klaidos.}
Skersmuo nėra tas pats, kas spindulys: \(d=2r\). Ploto formulėje kvadratu keliame tik spindulį, ne skaičių \(\pi\). Ilgis matuojamas ilgio vienetais, o plotas -- kvadratiniais vienetais.

\subsection{Iliustruojantys pavyzdžiai}
\textbf{1 pavyzdys.} Apskritimo spindulys \(5\) cm. Raskime skersmenį.
\[
d=2r=2\cdot 5=10\text{ cm}.
\]
\textbf{Atsakymas:} \(10\) cm.

\textbf{2 pavyzdys.} Raskime apskritimo ilgį, kai \(r=4\) cm.
\[
C=2\pi r=2\cdot 3,14\cdot 4=25,12\text{ cm}.
\]
\textbf{Atsakymas:} apie \(25,12\) cm.

\textbf{3 pavyzdys.} Raskime skritulio plotą, kai \(d=12\) m.
\[
r=\frac{12}{2}=6\text{ m},\qquad S=\pi r^2\approx 3,14\cdot 36=113,04\text{ m}^2.
\]
\textbf{Atsakymas:} apie \(113,04\text{ m}^2\).

\paragraph{Pasitikrinkite.}
\begin{enumerate}[label=\arabic*.]
\item Raskite skersmenį, kai \(r=7\) cm.
\item Raskite apskritimo ilgį, kai \(d=20\) cm.
\item Raskite skritulio plotą, kai \(r=3\) m.
\end{enumerate}

\subsection{Ryšys su kitomis temomis}
Apskritimo ir skritulio formulės naudoja daugybą, laipsnį \(r^2\), matavimo vienetus ir apytikslius skaičiavimus. Jos siejasi su plotu \ref{sec:plotas-pavirsius-ir-turis}, matavimais \ref{sec:matavimo-vienetu-keitimas-ir-kelio-formule} ir vėliau su geometrijos įrodymais.

\section{Plotas, paviršius ir tūris}\label{sec:plotas-pavirsius-ir-turis}

\subsection{Apibrėžtys}
\begin{apibreztis}
Perimetras yra plokščios figūros kraštinių ilgių suma.
\end{apibreztis}
\begin{apibreztis}
Plotas rodo, kiek kvadratinių vienetų telpa figūros viduje.
\end{apibreztis}
\begin{apibreztis}
Paviršiaus plotas yra visų erdvės kūno sienų plotų suma.
\end{apibreztis}
\begin{apibreztis}
Tūris rodo, kiek kubinių vienetų telpa erdvės kūno viduje.
\end{apibreztis}

\subsection{Teoremos ir savybės}
\begin{savybe}
Stačiakampio plotas \(S=ab\), perimetras \(P=2(a+b)\). Kvadrato plotas \(S=a^2\), perimetras \(P=4a\).
\end{savybe}
\begin{savybe}
Stačiojo trikampio plotas lygus pusei stačiakampio su tokiomis pačiomis statinėmis ploto:
\[
S=\frac{ab}{2}.
\]
\end{savybe}
\begin{savybe}
Stačiakampio gretasienio tūris
\[
V=abc,
\]
o kubo tūris \(V=a^3\).
\end{savybe}
\begin{savybe}
Sudėtinės figūros plotą galima rasti suskaidžius ją į žinomas figūras arba iš didesnės figūros atėmus trūkstamas dalis.
\end{savybe}

\subsection{Taisyklės ir algoritmai}
\textbf{Sudėtinės figūros plotas.}
\begin{enumerate}[label=\arabic*.]
\item Pažymėkite figūrą, kurią mokate apskaičiuoti: stačiakampį, kvadratą, stačiąjį trikampį.
\item Suskaidykite sudėtinę figūrą į šias dalis arba papildykite iki didesnės figūros.
\item Apskaičiuokite kiekvienos dalies plotą.
\item Sudėkite dalis arba atimkite nereikalingą plotą.
\item Užrašykite kvadratinius vienetus.
\end{enumerate}

\textbf{Stačiakampio gretasienio paviršius.}
\begin{enumerate}[label=\arabic*.]
\item Raskite trijų skirtingų sienų plotus: \(ab\), \(bc\), \(ac\).
\item Kiekviena tokia siena turi priešingą lygią sieną.
\item Skaičiuokite \(S_{\text{pav}}=2ab+2bc+2ac\).
\end{enumerate}

\paragraph{Dažnos klaidos.}
Perimetras ir plotas nėra tas pats: perimetras matuojamas cm, m, o plotas -- \(\text{cm}^2\), \(\text{m}^2\). Tūriui reikia trijų matmenų ir kubinių vienetų. Sudėtinėje figūroje nepadvigubinkite bendros vidinės kraštinės skaičiuodami perimetrą.

\subsection{Iliustruojantys pavyzdžiai}
\textbf{1 pavyzdys.} Stačiakampio ilgis \(8\) cm, plotis \(5\) cm. Raskime plotą ir perimetrą.
\[
S=8\cdot 5=40\text{ cm}^2,\qquad P=2(8+5)=26\text{ cm}.
\]
\textbf{Atsakymas:} \(40\text{ cm}^2\), \(26\) cm.

\textbf{2 pavyzdys.} Raskime stačiojo trikampio plotą, jei statinės \(6\) cm ir \(9\) cm.
\[
S=\frac{6\cdot 9}{2}=27\text{ cm}^2.
\]
\textbf{Atsakymas:} \(27\text{ cm}^2\).

\textbf{3 pavyzdys.} Stačiakampio gretasienio matmenys \(4\) cm, \(3\) cm ir \(10\) cm. Raskime tūrį.
\[
V=4\cdot 3\cdot 10=120\text{ cm}^3.
\]
\textbf{Atsakymas:} \(120\text{ cm}^3\).

\textbf{4 pavyzdys.} Kubo briauna \(5\) cm. Raskime paviršiaus plotą.
\[
S_{\text{pav}}=6a^2=6\cdot 25=150\text{ cm}^2.
\]
\textbf{Atsakymas:} \(150\text{ cm}^2\).

\paragraph{Pasitikrinkite.}
\begin{enumerate}[label=\arabic*.]
\item Kvadrato kraštinė \(11\) cm. Raskite perimetrą ir plotą.
\item Stačiojo trikampio statinės \(8\) m ir \(5\) m. Raskite plotą.
\item Stačiakampio gretasienio matmenys \(2\), \(6\), \(7\) cm. Raskite tūrį.
\end{enumerate}

\subsection{Ryšys su kitomis temomis}
Plotai, paviršiai ir tūriai jungia geometriją su matavimo vienetais \ref{sec:matavimo-vienetu-keitimas-ir-kelio-formule}, trupmenomis, dešimtainiais skaičiais ir lygtimis. Jie reikalingi praktiniams statybos, pakavimo, žemėlapių ir modelių uždaviniams.

\section{Statistiniai duomenys ir vidurkiai}\label{sec:statistiniai-duomenys-ir-vidurkiai}

\subsection{Apibrėžtys}
\begin{apibreztis}
Duomenys yra surinktos reikšmės apie tiriamą požymį, pavyzdžiui, mokinių ūgiai arba mėgstamos sporto šakos.
\end{apibreztis}
\begin{apibreztis}
Populiacija yra visa tiriamų objektų grupė, o imtis -- jos dalis, iš kurios renkami duomenys.
\end{apibreztis}
\begin{apibreztis}
Imties vidurkis gaunamas sudėjus visas kiekybines reikšmes ir padalijus iš reikšmių skaičiaus.
\end{apibreztis}
\begin{apibreztis}
Moda yra dažniausiai pasikartojanti reikšmė, o mediana yra vidurinė surikiuotų duomenų reikšmė.
\end{apibreztis}

\subsection{Teoremos ir savybės}
\begin{savybe}
Vidurkį prasminga skaičiuoti kiekybiniams duomenims, pavyzdžiui, laikui, ilgiui, taškams, kainoms.
\end{savybe}
\begin{savybe}
Mediana mažiau priklauso nuo labai didelių arba labai mažų išskirtinių reikšmių negu vidurkis.
\end{savybe}
\begin{savybe}
Dažnių lentelėje reikšmės dažnis parodo, kiek kartų ši reikšmė pasikartojo.
\end{savybe}
\begin{savybe}
Diagramos turi turėti aiškų pavadinimą, ašių arba kategorijų pavadinimus ir vienetus, kitaip jas galima klaidingai interpretuoti.
\end{savybe}

\subsection{Taisyklės ir algoritmai}
\textbf{Statistinis tyrimas.}
\begin{enumerate}[label=\arabic*.]
\item Suformuluokite klausimą, į kurį galima atsakyti duomenimis.
\item Nuspręskite, kokią imtį tirsime.
\item Surinkite duomenis tuo pačiu būdu iš visų imties narių.
\item Sutvarkykite duomenis lentelėje.
\item Apskaičiuokite tinkamas charakteristikas ir padarykite išvadą.
\end{enumerate}

\textbf{Medianos radimas.}
\begin{enumerate}[label=\arabic*.]
\item Surikiuokite duomenis nuo mažiausio iki didžiausio.
\item Jei reikšmių skaičius nelyginis, mediana yra vidurinė reikšmė.
\item Jei reikšmių skaičius lyginis, mediana yra dviejų vidurinių reikšmių vidurkis.
\end{enumerate}

\paragraph{Dažnos klaidos.}
Vidurkis nėra būtinai viena iš duomenų reikšmių. Kokybiniams duomenims, pavyzdžiui, spalvoms, vidurkio skaičiuoti negalima. Diagramose reikia žiūrėti į skalę: vienas padalos tarpas gali reikšti ne \(1\), o \(5\), \(10\) ar \(100\) vienetų.

\subsection{Iliustruojantys pavyzdžiai}
\textbf{1 pavyzdys.} Penki mokiniai perskaitė \(3\), \(5\), \(5\), \(7\), \(10\) knygų. Raskime vidurkį, modą ir medianą.
\[
\overline{x}=\frac{3+5+5+7+10}{5}=\frac{30}{5}=6.
\]
Dažniausiai kartojasi \(5\), todėl moda \(5\). Surikiuota eilė jau duota, vidurinė reikšmė \(5\), todėl mediana \(5\).

\textbf{Atsakymas:} vidurkis \(6\), moda \(5\), mediana \(5\).

\textbf{2 pavyzdys.} Dažnių lentelėje pateikti kontrolinio darbo pažymiai.
\[
\begin{array}{c|cccc}
\text{Pažymys}&7&8&9&10\\ \hline
\text{Dažnis}&3&5&4&2
\end{array}
\]
Kiek mokinių rašė darbą?
\[
3+5+4+2=14.
\]
\textbf{Atsakymas:} \(14\) mokinių.

\textbf{3 pavyzdys.} Raskime šių duomenų medianą: \(12,9,15,10\).

Rikiuojame:
\[
9,10,12,15.
\]
Vidurinės reikšmės yra \(10\) ir \(12\), todėl
\[
\text{mediana}=\frac{10+12}{2}=11.
\]
\textbf{Atsakymas:} \(11\).

\paragraph{Pasitikrinkite.}
\begin{enumerate}[label=\arabic*.]
\item Raskite duomenų \(4,6,6,8,11\) vidurkį, modą ir medianą.
\item Ar galima skaičiuoti mėgstamiausių spalvų vidurkį?
\item Dažniai yra \(2,7,3,1\). Kiek iš viso duomenų?
\end{enumerate}

\subsection{Ryšys su kitomis temomis}
Statistika naudoja skaičiavimą, lenteles, diagramas ir proporcijas. Ji siejasi su tikimybe \ref{sec:atsitiktiniai-bandymas-ir-tikimybe}, nes ilgai kartojant atsitiktinius bandymus santykiniai dažniai padeda vertinti įvykio tikėtinumą.

\section{Atsitiktinis bandymas ir tikimybė}\label{sec:atsitiktiniai-bandymas-ir-tikimybe}

\subsection{Apibrėžtys}
\begin{apibreztis}
Atsitiktinis bandymas yra veiksmas arba situacija, kurios tikslios baigties iš anksto nežinome, bet galime išvardyti galimas baigtis.
\end{apibreztis}
\begin{apibreztis}
Baigtis yra vienas galimas bandymo rezultatas. Baigčių aibė yra visų galimų baigčių rinkinys.
\end{apibreztis}
\begin{apibreztis}
Įvykis yra viena baigtis arba baigčių rinkinys. Elementarusis įvykis turi vieną baigtį, būtinasis įvykis įvyksta visada, negalimasis -- niekada.
\end{apibreztis}
\begin{apibreztis}
Kai visos baigtys vienodai galimos, įvykio tikimybė yra palankių baigčių skaičiaus ir visų baigčių skaičiaus santykis:
\[
P(A)=\frac{\text{palankių baigčių skaičius}}{\text{visų baigčių skaičius}}.
\]
\end{apibreztis}

\subsection{Teoremos ir savybės}
\begin{savybe}
Kiekvieno įvykio tikimybė yra skaičius nuo \(0\) iki \(1\): \(0\le P(A)\le 1\).
\end{savybe}
\begin{savybe}
Būtinojo įvykio tikimybė lygi \(1\), negalimojo įvykio tikimybė lygi \(0\).
\end{savybe}
\begin{savybe}
Įvykio ir jam priešingo įvykio tikimybių suma lygi \(1\):
\[
P(A)+P(\text{ne }A)=1.
\]
\end{savybe}
\begin{savybe}
Jei pirmame etape yra \(m\) galimybių, o po kiekvienos jų antrame etape yra \(n\) galimybių, tai dviejų etapų baigčių yra \(m\cdot n\).
\end{savybe}

\subsection{Taisyklės ir algoritmai}
\textbf{Tikimybės skaičiavimas vienodų baigčių atveju.}
\begin{enumerate}[label=\arabic*.]
\item Išvardykite visas galimas baigtis.
\item Patikrinkite, ar jos vienodai galimos.
\item Suskaičiuokite, kiek baigčių palankios įvykiui.
\item Taikykite formulę \(P(A)=\frac{m}{n}\).
\item Atsakymą galite užrašyti trupmena, dešimtaine trupmena arba procentais.
\end{enumerate}

\textbf{Dviejų etapų bandymas.}
\begin{enumerate}[label=\arabic*.]
\item Sudarykite galimybių medį arba lentelę.
\item Kiekvieną kelią nuo pradžios iki pabaigos laikykite viena baigtimi.
\item Suskaičiuokite visas ir palankias baigtis.
\end{enumerate}

\paragraph{Dažnos klaidos.}
Ne visos baigtys visada vienodai galimos. Pavyzdžiui, „laimės komanda A“ ir „laimės komanda B“ nebūtinai turi vienodą tikimybę. Tikimybė negali būti didesnė už \(1\) arba mažesnė už \(0\). Priešingas įvykis reiškia „neįvyks“, o ne „įvyks priešingai pagal nuotaiką“.

\subsection{Iliustruojantys pavyzdžiai}
\textbf{1 pavyzdys.} Metamas sąžiningas šešiasienis kauliukas. Kokia tikimybė išridenti lyginį skaičių?

Baigtys: \(1,2,3,4,5,6\). Palankios: \(2,4,6\), jų yra \(3\).
\[
P=\frac{3}{6}=\frac{1}{2}.
\]
\textbf{Atsakymas:} \(\frac{1}{2}\).

\textbf{2 pavyzdys.} Maišelyje yra \(3\) raudoni ir \(2\) mėlyni rutuliukai. Atsitiktinai traukiamas vienas rutuliukas. Kokia tikimybė ištraukti mėlyną?
\[
P(\text{mėlynas})=\frac{2}{3+2}=\frac{2}{5}.
\]
\textbf{Atsakymas:} \(\frac{2}{5}\).

\textbf{3 pavyzdys.} Metama moneta ir sukamas suktukas su skaičiais \(1,2,3\). Kiek yra baigčių?

Moneta turi \(2\) baigtis, suktukas \(3\), todėl
\[
2\cdot 3=6.
\]
Baigtys gali būti \((H;1),(H;2),(H;3),(S;1),(S;2),(S;3)\).

\textbf{Atsakymas:} \(6\) baigtys.

\paragraph{Pasitikrinkite.}
\begin{enumerate}[label=\arabic*.]
\item Kokia tikimybė išridenti skaičių, didesnį už \(4\), metant kauliuką?
\item Maišelyje \(5\) balti ir \(5\) juodi rutuliukai. Kokia tikimybė ištraukti baltą?
\item Kiek baigčių yra renkantis vieną iš \(4\) marškinėlių ir vieną iš \(3\) kepurių?
\end{enumerate}

\subsection{Ryšys su kitomis temomis}
Tikimybė remiasi trupmenomis, procentais ir duomenų analizės idėjomis. Ji siejasi su statistika \ref{sec:statistiniai-duomenys-ir-vidurkiai}, nes eksperimentais galima apytiksliai įvertinti atsitiktinių įvykių tikėtinumą.

\section{Romėniškieji skaičiai ir pozicinė sistema}\label{sec:romeniskieji-skaiciai-ir-pozicine-sistema}

\subsection{Apibrėžtys}
\begin{apibreztis}
Romėniškieji skaitmenys yra \(I=1\), \(V=5\), \(X=10\), \(L=50\), \(C=100\), \(D=500\), \(M=1000\).
\end{apibreztis}
\begin{apibreztis}
Dešimtainė pozicinė sistema yra skaičių rašymo sistema, kurioje skaitmens reikšmė priklauso nuo jo vietos skaičiuje.
\end{apibreztis}
\begin{apibreztis}
Skyriaus reikšmė yra skaitmens indėlis į skaičių: skaičiuje \(4062\) skaitmuo \(4\) reiškia \(4000\), o \(6\) reiškia \(60\).
\end{apibreztis}

\subsection{Teoremos ir savybės}
\begin{savybe}
Dešimtainėje sistemoje kiekviena vieta į kairę yra \(10\) kartų didesnė, o kiekviena vieta į dešinę po kablelio yra \(10\) kartų mažesnė.
\end{savybe}
\begin{savybe}
Romėniškuose skaičiuose mažesnis ženklas prieš didesnį dažniausiai atimamas: \(IV=4\), \(IX=9\), \(XL=40\), \(XC=90\), \(CD=400\), \(CM=900\).
\end{savybe}
\begin{savybe}
Romėniškoji sistema nėra pozicinė: ženklas \(X\) visada reiškia \(10\), o ne priklausomai nuo vietos kintančią reikšmę.
\end{savybe}

\subsection{Taisyklės ir algoritmai}
\textbf{Natūraliojo skaičiaus rašymas romėniškai.}
\begin{enumerate}[label=\arabic*.]
\item Išskaidykite skaičių į tūkstančius, šimtus, dešimtis ir vienetus.
\item Kiekvieną dalį užrašykite romėniškais ženklais.
\item Sujunkite dalis nuo didžiausios iki mažiausios.
\end{enumerate}

\textbf{Skaičiaus skaidymas skyrių suma.}
\begin{enumerate}[label=\arabic*.]
\item Perskaitykite skaitmenų vietas.
\item Kiekvieną nenulinį skaitmenį padauginkite iš jo skyriaus vieneto.
\item Sudėkite gautas dalis.
\end{enumerate}

\paragraph{Dažnos klaidos.}
Romėniškai nerašome keturių vienodų ženklų iš eilės: \(4\) yra \(IV\), ne \(IIII\). Nulis romėniškoje sistemoje įprastai nenaudojamas. Dešimtainėje sistemoje nulis gali būti labai svarbus, nes palaiko vietą: \(305\ne 35\).

\subsection{Iliustruojantys pavyzdžiai}
\textbf{1 pavyzdys.} Užrašykime \(1987\) romėniškai.
\[
1987=1000+900+80+7.
\]
\[
1000=M,\quad 900=CM,\quad 80=LXXX,\quad 7=VII.
\]
\textbf{Atsakymas:} \(MCMLXXXVII\).

\textbf{2 pavyzdys.} Perskaitykime \(XLIV\).
\[
XL=40,\qquad IV=4,
\]
todėl \(XLIV=44\).

\textbf{Atsakymas:} \(44\).

\textbf{3 pavyzdys.} Išskaidykime \(7\,305\,040\) skyrių suma.
\[
7\,305\,040=7\,000\,000+300\,000+5\,000+40.
\]
\textbf{Atsakymas:} \(7\,000\,000+300\,000+5\,000+40\).

\paragraph{Pasitikrinkite.}
\begin{enumerate}[label=\arabic*.]
\item Užrašykite \(246\) romėniškai.
\item Perskaitykite \(MMXXVI\).
\item Išskaidykite \(50\,608\) skyrių suma.
\end{enumerate}

\subsection{Ryšys su kitomis temomis}
Pozicinė sistema padeda suprasti apvalinimą, dešimtaines trupmenas ir veiksmus stulpeliu. Romėniškieji skaičiai parodo, kad skaičių rašymo būdai gali būti skirtingi, todėl verta atskirti skaičių nuo jo užrašo.

\section{Dalumo požymiai, pirminiai skaičiai, BDD ir MBK}\label{sec:dalumo-pozymiai-pirminiai-skaiciai-bdd-mbk}

\subsection{Apibrėžtys}
\begin{apibreztis}
Bendrasis daliklis yra skaičius, iš kurio dalijasi du ar daugiau skaičių. Didžiausias bendrasis daliklis, trumpai DBD, yra didžiausias iš bendrųjų daliklių.
\end{apibreztis}
\begin{apibreztis}
Bendrasis kartotinis yra skaičius, kuris dalijasi iš duotų skaičių. Mažiausias teigiamas bendrasis kartotinis, trumpai MBK, yra mažiausias iš teigiamų bendrųjų kartotinių.
\end{apibreztis}
\begin{apibreztis}
Pirminis skaidymas yra skaičiaus užrašymas pirminių skaičių sandauga.
\end{apibreztis}

\subsection{Teoremos ir savybės}
\begin{savybe}
Jei du skaičiai neturi bendrų daliklių, išskyrus \(1\), jų DBD lygus \(1\). Tokie skaičiai vadinami tarpusavyje pirminiais.
\end{savybe}
\begin{savybe}
Jei \(a\mid b\), tai \(\operatorname{DBD}(a,b)=a\), o \(\operatorname{MBK}(a,b)=b\).
\end{savybe}
\begin{savybe}
Teigiamiems skaičiams \(a\) ir \(b\) galioja
\[
\operatorname{DBD}(a,b)\cdot \operatorname{MBK}(a,b)=a\cdot b.
\]
\end{savybe}

\subsection{Taisyklės ir algoritmai}
\textbf{Dalumo požymių naudojimas.}
\begin{enumerate}[label=\arabic*.]
\item Iš \(2\), \(5\), \(10\) tikrinkite pagal paskutinį skaitmenį.
\item Iš \(3\) ir \(9\) tikrinkite pagal skaitmenų sumą.
\item Iš \(4\) tikrinkite, ar iš \(4\) dalijasi paskutiniai du skaitmenys.
\item Iš \(6\) skaičius dalijasi tada, kai dalijasi ir iš \(2\), ir iš \(3\).
\end{enumerate}

\textbf{Pirminio skaidymo medis.}
\begin{enumerate}[label=\arabic*.]
\item Skaičių skaidykite į dviejų natūraliųjų skaičių sandaugą.
\item Kiekvieną sudėtinį dauginamąjį skaidykite toliau.
\item Baikite, kai visi šakų galai yra pirminiai skaičiai.
\item Gautus pirminius dauginamuosius užrašykite sandauga.
\end{enumerate}

\paragraph{Dažnos klaidos.}
DBD ieško bendrų daliklių, MBK -- bendrų kartotinių. Pirminio skaidymo tvarka gali skirtis, bet galutiniai pirminiai dauginamieji turi sutapti. Dalumo iš \(6\) požymiui neužtenka, kad skaičius dalytųsi tik iš \(2\) arba tik iš \(3\).

\subsection{Iliustruojantys pavyzdžiai}
\textbf{1 pavyzdys.} Patikrinkime, ar \(2316\) dalijasi iš \(4\), \(6\), \(9\).

Paskutiniai du skaitmenys \(16\) dalijasi iš \(4\), todėl \(2316\) dalijasi iš \(4\). Skaičius lyginis, o skaitmenų suma
\[
2+3+1+6=12
\]
dalijasi iš \(3\), todėl skaičius dalijasi iš \(6\). Kadangi \(12\) nesidalija iš \(9\), skaičius nesidalija iš \(9\).

\textbf{Atsakymas:} dalijasi iš \(4\) ir \(6\), nesidalija iš \(9\).

\textbf{2 pavyzdys.} Raskime \(\operatorname{DBD}(48,180)\).
\[
48=2^4\cdot 3,\qquad 180=2^2\cdot 3^2\cdot 5.
\]
Bendri mažiausi laipsniai: \(2^2\cdot 3\), todėl
\[
\operatorname{DBD}(48,180)=12.
\]
\textbf{Atsakymas:} \(12\).

\textbf{3 pavyzdys.} Raskime \(\operatorname{MBK}(15,18,20)\).
\[
15=3\cdot 5,\quad 18=2\cdot 3^2,\quad 20=2^2\cdot 5.
\]
Imame didžiausius laipsnius:
\[
\operatorname{MBK}=2^2\cdot 3^2\cdot 5=180.
\]
\textbf{Atsakymas:} \(180\).

\paragraph{Pasitikrinkite.}
\begin{enumerate}[label=\arabic*.]
\item Ar \(3276\) dalijasi iš \(4\), \(6\), \(9\)?
\item Raskite \(\operatorname{DBD}(72,96)\).
\item Raskite \(\operatorname{MBK}(12,15,25)\).
\end{enumerate}

\subsection{Ryšys su kitomis temomis}
Ši tema gilina \ref{sec:naturalieji-skaiciai-dalikliai-ir-kartotiniai}. DBD reikalingas trupmenoms trumpinti, MBK -- bendram vardikliui rasti, o dalumo požymiai padeda skaičiuoti greičiau ir tikrinti atsakymus.

\section{Trupmenų veiksmai ir procentinis pokytis}\label{sec:trupmenu-veiksmai-ir-procentinis-pokytis}

\subsection{Apibrėžtys}
\begin{apibreztis}
Mišrusis skaičius yra skaičius, sudarytas iš sveikosios dalies ir taisyklingosios trupmenos, pavyzdžiui, \(2\frac{3}{5}\).
\end{apibreztis}
\begin{apibreztis}
Skaičiui atvirkštinis skaičius yra toks skaičius, kurio sandauga su pradiniu skaičiumi lygi \(1\). Skaičiui \(\frac{a}{b}\) atvirkštinis yra \(\frac{b}{a}\), kai \(a\ne 0\).
\end{apibreztis}
\begin{apibreztis}
Procentinis padidėjimas reiškia, kad prie pradinės reikšmės pridedama nurodyta jos procentinė dalis. Procentinis sumažėjimas reiškia, kad ši dalis atimama.
\end{apibreztis}

\subsection{Teoremos ir savybės}
\begin{savybe}
Trupmenų daugyba:
\[
\frac{a}{b}\cdot \frac{c}{d}=\frac{ac}{bd},\qquad b\ne 0,\ d\ne 0.
\]
\end{savybe}
\begin{savybe}
Dalyba iš trupmenos pakeičiama daugyba iš atvirkštinės trupmenos:
\[
\frac{a}{b}:\frac{c}{d}=\frac{a}{b}\cdot \frac{d}{c},\qquad c\ne 0.
\]
\end{savybe}
\begin{savybe}
Padidinus skaičių \(p\%\), nauja reikšmė lygi \(a(1+\frac{p}{100})\). Sumažinus \(p\%\), nauja reikšmė lygi \(a(1-\frac{p}{100})\).
\end{savybe}

\subsection{Taisyklės ir algoritmai}
\textbf{Mišriųjų skaičių veiksmai.}
\begin{enumerate}[label=\arabic*.]
\item Sudėčiai ir atimčiai galima atskirai tvarkyti sveikąsias ir trupmenines dalis, bet jei trupmeninė dalis „persiskolina“, patogu mišriuosius skaičius paversti netaisyklingosiomis trupmenomis.
\item Daugybai ir dalybai mišriuosius skaičius visada paverskite netaisyklingosiomis trupmenomis.
\item Atlikite veiksmą ir sutrumpinkite rezultatą.
\item Jei reikia, rezultatą užrašykite mišriuoju skaičiumi.
\end{enumerate}

\textbf{Procentinio pokyčio uždavinys.}
\begin{enumerate}[label=\arabic*.]
\item Nustatykite pradinę reikšmę.
\item Procentą pakeiskite dešimtaine trupmena.
\item Raskite pokyčio dydį arba naudokite daugiklį \(1\pm \frac{p}{100}\).
\item Padidėjimo atveju pridėkite, sumažėjimo atveju atimkite.
\end{enumerate}

\paragraph{Dažnos klaidos.}
Dalyba iš trupmenos nėra dalyba iš skaitiklio ir vardiklio atskirai. Reikia dauginti iš atvirkštinės trupmenos. Procentinis padidėjimas \(20\%\), po to sumažėjimas \(20\%\), negrąžina pradinės reikšmės, nes antras pokytis skaičiuojamas nuo naujos reikšmės.

\subsection{Iliustruojantys pavyzdžiai}
\textbf{1 pavyzdys.} Apskaičiuokime \(2\frac{1}{3}+1\frac{5}{6}\).
\[
2\frac{1}{3}=2\frac{2}{6},\qquad 1\frac{5}{6}=1\frac{5}{6}.
\]
\[
2\frac{2}{6}+1\frac{5}{6}=3\frac{7}{6}=4\frac{1}{6}.
\]
\textbf{Atsakymas:} \(4\frac{1}{6}\).

\textbf{2 pavyzdys.} Apskaičiuokime \(\frac{5}{8}:\frac{15}{16}\).
\[
\frac{5}{8}:\frac{15}{16}=\frac{5}{8}\cdot \frac{16}{15}.
\]
Trumpiname:
\[
\frac{5\cdot 16}{8\cdot 15}=\frac{1\cdot 2}{1\cdot 3}=\frac{2}{3}.
\]
\textbf{Atsakymas:} \(\frac{2}{3}\).

\textbf{3 pavyzdys.} Skaičių \(240\) padidinkime \(15\%\).
\[
240\cdot 1,15=276.
\]
\textbf{Atsakymas:} \(276\).

\textbf{4 pavyzdys.} Kaina \(60\) Eur sumažėjo \(25\%\). Raskime naują kainą.
\[
60\cdot 0,75=45.
\]
\textbf{Atsakymas:} \(45\) Eur.

\paragraph{Pasitikrinkite.}
\begin{enumerate}[label=\arabic*.]
\item Apskaičiuokite \(1\frac{3}{4}+2\frac{1}{2}\).
\item Apskaičiuokite \(\frac{7}{9}\cdot \frac{3}{14}\).
\item Skaičių \(80\) sumažinkite \(30\%\).
\end{enumerate}

\subsection{Ryšys su kitomis temomis}
Trupmenų veiksmai remiasi \ref{sec:paprastosios-ir-desimtaines-trupmenos} ir DBD bei MBK iš \ref{sec:dalumo-pozymiai-pirminiai-skaiciai-bdd-mbk}. Procentinis pokytis tiesiogiai taikomas finansiniuose skaičiavimuose \ref{sec:finansai-biudzetas-nuolaidos-ir-palukanos}.

\section{Finansiniai skaičiavimai: biudžetas, nuolaidos ir palūkanos}\label{sec:finansai-biudzetas-nuolaidos-ir-palukanos}

\subsection{Apibrėžtys}
\begin{apibreztis}
Biudžetas yra planas, kiek pinigų gaunama, kiek išleidžiama ir kiek lieka.
\end{apibreztis}
\begin{apibreztis}
Nuolaida yra kainos sumažinimas. Procentinė nuolaida nurodo, kokia pradinės kainos dalis atimama.
\end{apibreztis}
\begin{apibreztis}
Vieneto kaina arba tarifas rodo, kiek kainuoja vienas matavimo vienetas: \(1\) kg, \(1\) l, \(1\) kWh, \(1\) bilietas ir panašiai.
\end{apibreztis}
\begin{apibreztis}
Paprastosios palūkanos yra atlygis už pasiskolintus arba padėtus pinigus, skaičiuojamas nuo pradinės sumos.
\end{apibreztis}

\subsection{Teoremos ir savybės}
\begin{savybe}
Po \(p\%\) nuolaidos lieka \(100\%-p\%\) pradinės kainos.
\end{savybe}
\begin{savybe}
Vieneto kaina apskaičiuojama bendrą kainą dalijant iš kiekio:
\[
\text{vieneto kaina}=\frac{\text{bendra kaina}}{\text{kiekis}}.
\]
\end{savybe}
\begin{savybe}
Paprastosios palūkanos per vienus metus gali būti skaičiuojamos formule
\[
I=\frac{p}{100}\cdot S,
\]
kur \(S\) -- pradinė suma, \(p\) -- metinių palūkanų procentas.
\end{savybe}
\begin{savybe}
Finansinis sprendimas turi būti grindžiamas ne tik mažiausia kaina, bet ir kiekiu, kokybe, poreikiu bei turimu biudžetu.
\end{savybe}

\subsection{Taisyklės ir algoritmai}
\textbf{Biudžeto tikrinimas.}
\begin{enumerate}[label=\arabic*.]
\item Sudėkite visas pajamas.
\item Sudėkite visas planuojamas išlaidas.
\item Apskaičiuokite likutį: pajamos minus išlaidos.
\item Jei likutis neigiamas, reikia mažinti išlaidas arba ieškoti papildomų pajamų.
\end{enumerate}

\textbf{Pirkinių palyginimas.}
\begin{enumerate}[label=\arabic*.]
\item Jei kiekiai skirtingi, raskite vieneto kainą.
\item Jei taikoma nuolaida, pirmiausia apskaičiuokite galutinę kainą.
\item Lyginkite tos pačios rūšies vienetus.
\item Atsakymą pagrįskite skaičiavimu.
\end{enumerate}

\paragraph{Dažnos klaidos.}
Nuolaida skaičiuojama nuo pradinės kainos, nebent sąlygoje nurodyta kitaip. Pakuotė su didesniu kiekiu ne visada pigesnė už vienetą, todėl reikia skaičiuoti vieneto kainą. Planuojant biudžetą negalima pamiršti mažų pasikartojančių išlaidų.

\subsection{Iliustruojantys pavyzdžiai}
\textbf{1 pavyzdys.} Kuprinė kainavo \(48\) Eur, jai taikoma \(25\%\) nuolaida. Kokia galutinė kaina?
\[
48\cdot 0,75=36.
\]
\textbf{Atsakymas:} \(36\) Eur.

\textbf{2 pavyzdys.} Kuris jogurtas pigesnis: \(400\) g už \(1,60\) Eur ar \(250\) g už \(1,10\) Eur?

Pirmojo kaina už \(100\) g:
\[
1,60:4=0,40\text{ Eur}.
\]
Antrojo kaina už \(100\) g:
\[
1,10:2,5=0,44\text{ Eur}.
\]
\textbf{Atsakymas:} pigesnis pirmasis jogurtas.

\textbf{3 pavyzdys.} Mokinys per savaitę gauna \(12\) Eur. Jis planuoja išleisti \(4,50\) Eur už pietus, \(3,20\) Eur už transportą ir \(2,80\) Eur už pramogas. Kiek liks?
\[
4,50+3,20+2,80=10,50,\qquad 12-10,50=1,50.
\]
\textbf{Atsakymas:} liks \(1,50\) Eur.

\textbf{4 pavyzdys.} Indėlis \(200\) Eur vieniems metams duoda \(4\%\) paprastųjų palūkanų. Kiek palūkanų gaunama?
\[
I=0,04\cdot 200=8.
\]
\textbf{Atsakymas:} \(8\) Eur.

\paragraph{Pasitikrinkite.}
\begin{enumerate}[label=\arabic*.]
\item Prekė kainuoja \(35\) Eur, nuolaida \(20\%\). Raskite galutinę kainą.
\item \(2\) kg ryžių kainuoja \(3,80\) Eur. Kiek kainuoja \(1\) kg?
\item Pajamos \(18\) Eur, išlaidos \(6,40\), \(5,90\), \(2,10\) Eur. Koks likutis?
\end{enumerate}

\subsection{Ryšys su kitomis temomis}
Finansiniai skaičiavimai taiko procentus \ref{sec:procentai-santykiai-ir-proporcijos-pradzia}, dešimtainius skaičius, santykius ir paprastas lygtis. Tai viena aiškiausių temų, kur matematika padeda priimti kasdienius sprendimus.

\section{Tiesiniai ir netiesiniai sąryšiai}\label{sec:tiesiniai-ir-netiesiniai-sarysiai-5-6}

\subsection{Apibrėžtys}
\begin{apibreztis}
Sąryšis nusako, kaip vienas dydis susijęs su kitu. Sąryšį galima pateikti žodžiais, lentele, formule, skaičių poromis arba grafiku.
\end{apibreztis}
\begin{apibreztis}
Tiesinis sąryšis yra toks, kurio grafikas koordinačių plokštumoje yra tiesė arba tiesės dalis. Paprasti pavyzdžiai: \(y=kx\), \(y=kx+b\).
\end{apibreztis}
\begin{apibreztis}
Tiesioginis proporcingumas yra tiesinis sąryšis \(y=kx\), kai santykis \(\frac{y}{x}\) yra pastovus visoms \(x\ne 0\) reikšmėms.
\end{apibreztis}
\begin{apibreztis}
Netiesinis sąryšis nėra vaizduojamas viena tiese. Pavyzdžiui, kvadrato plotas \(S=a^2\) nuo kraštinės \(a\) priklauso netiesiškai.
\end{apibreztis}

\subsection{Teoremos ir savybės}
\begin{savybe}
Tiesioginio proporcingumo grafikas eina per pradžios tašką \((0;0)\), jei nagrinėjamos visos tinkamos reikšmės.
\end{savybe}
\begin{savybe}
Jei lentelėje santykis \(\frac{y}{x}\) visur tas pats, dydžiai yra tiesiogiai proporcingi.
\end{savybe}
\begin{savybe}
Sąryšyje \(y=kx+b\) vienodai padidinus \(x\), \(y\) kinta vienodu dydžiu.
\end{savybe}
\begin{savybe}
Netiesiniame sąryšyje vienodai keičiant \(x\), \(y\) pokytis gali būti nevienodas.
\end{savybe}

\subsection{Taisyklės ir algoritmai}
\textbf{Ar sąryšis tiesiogiai proporcingas?}
\begin{enumerate}[label=\arabic*.]
\item Patikrinkite, ar \(x\ne 0\) reikšmėms galima skaičiuoti \(\frac{y}{x}\).
\item Apskaičiuokite šį santykį keliose lentelės eilutėse.
\item Jei santykis vienodas, sąryšis tiesiogiai proporcingas.
\item Užrašykite formulę \(y=kx\), kur \(k\) yra rastas santykis.
\end{enumerate}

\textbf{Grafiko skaitymas.}
\begin{enumerate}[label=\arabic*.]
\item Pasirinkite tašką grafike.
\item Nuo taško nusileiskite iki \(x\) ašies ir perskaitykite \(x\).
\item Nuo taško eikite iki \(y\) ašies ir perskaitykite \(y\).
\item Atsakymą aiškinkite pagal uždavinio kontekstą.
\end{enumerate}

\paragraph{Dažnos klaidos.}
Ne kiekviena didėjanti lentelė yra tiesioginis proporcingumas. Jei \(x=0\), formulėje \(y=kx\) turi būti \(y=0\). Jei lentelėje \(x\) didėja po \(1\), o \(y\) didėja po \(2\), sąryšis tiesinis, bet nebūtinai tiesiogiai proporcingas, pavyzdžiui, \(y=2x+5\).

\subsection{Iliustruojantys pavyzdžiai}
\textbf{1 pavyzdys.} Patikrinkime, ar lentelė rodo tiesioginį proporcingumą.
\[
\begin{array}{c|ccc}
x&2&5&8\\ \hline
y&6&15&24
\end{array}
\]
\[
\frac{6}{2}=3,\qquad \frac{15}{5}=3,\qquad \frac{24}{8}=3.
\]
Santykis pastovus, todėl \(y=3x\).

\textbf{Atsakymas:} taip, \(y=3x\).

\textbf{2 pavyzdys.} Lentelėje \(x=1,2,3,4\), o \(y=4,7,10,13\). Ar tai tiesioginis proporcingumas?

Santykiai:
\[
\frac{4}{1}=4,\qquad \frac{7}{2}=3,5.
\]
Jie nevienodi, todėl ne tiesioginis proporcingumas. Tačiau \(y\) didėja po \(3\), todėl sąryšis tiesinis: \(y=3x+1\).

\textbf{Atsakymas:} ne tiesioginis proporcingumas, bet tiesinis sąryšis.

\textbf{3 pavyzdys.} Kvadrato plotas \(S=a^2\). Raskime plotą, kai \(a=1,2,3\).
\[
1^2=1,\qquad 2^2=4,\qquad 3^2=9.
\]
Pokyčiai \(+3\) ir \(+5\) nevienodi, todėl sąryšis netiesinis.

\textbf{Atsakymas:} plotai \(1,4,9\), sąryšis netiesinis.

\paragraph{Pasitikrinkite.}
\begin{enumerate}[label=\arabic*.]
\item Ar lentelė \((2;10)\), \((3;15)\), \((6;30)\) rodo tiesioginį proporcingumą?
\item Raskite formulę, jei \(y\) visada \(4\) kartus didesnis už \(x\).
\item Ar kubo tūris \(V=a^3\) nuo briaunos \(a\) priklauso tiesiškai?
\end{enumerate}

\subsection{Ryšys su kitomis temomis}
Sąryšiai jungia proporcijas \ref{sec:procentai-santykiai-ir-proporcijos-pradzia}, koordinates \ref{sec:dydziu-priklausomybes-ir-koordinaciu-plokstuma}, lygtis ir praktinius judėjimo bei finansų uždavinius.

\section{Matavimo vienetų keitimas ir kelio formulė}\label{sec:matavimo-vienetu-keitimas-ir-kelio-formule}

\subsection{Apibrėžtys}
\begin{apibreztis}
Matavimo vienetas yra sutartas dydis, su kuriuo lyginame matuojamą objektą, pavyzdžiui, metras, sekundė, litras.
\end{apibreztis}
\begin{apibreztis}
Kelias \(s\) yra nueitas arba nuvažiuotas atstumas, greitis \(v\) rodo, koks kelias įveikiamas per laiko vienetą, laikas \(t\) rodo judėjimo trukmę.
\end{apibreztis}
\begin{apibreztis}
Kelio formulė yra
\[
s=vt.
\]
Iš jos gauname \(v=\frac{s}{t}\) ir \(t=\frac{s}{v}\), kai dalikliai nelygūs nuliui.
\end{apibreztis}

\subsection{Teoremos ir savybės}
\begin{savybe}
Metrinėje sistemoje ilgio vienetai keičiami dauginant arba dalijant iš \(10\), \(100\), \(1000\) ir panašiai: \(1\text{ km}=1000\text{ m}\), \(1\text{ m}=100\text{ cm}\).
\end{savybe}
\begin{savybe}
Ploto vienetai keičiasi kvadratiškai: \(1\text{ m}^2=10\,000\text{ cm}^2\). Tūrio vienetai keičiasi kubiškai: \(1\text{ m}^3=1\,000\,000\text{ cm}^3\).
\end{savybe}
\begin{savybe}
Jei du kūnai juda vienas prieš kitą, jų artėjimo greitis yra greičių suma. Jei juda ta pačia kryptimi, tolimo arba vijimosi greitis dažnai yra greičių skirtumas.
\end{savybe}

\subsection{Taisyklės ir algoritmai}
\textbf{Vienetų keitimas.}
\begin{enumerate}[label=\arabic*.]
\item Užrašykite, iš kokio vieneto į kokį keičiate.
\item Nustatykite, ar naujas vienetas mažesnis, ar didesnis.
\item Į mažesnius vienetus skaitinė reikšmė didėja, į didesnius -- mažėja.
\item Patikrinkite, ar ploto ir tūrio vienetams pritaikėte kvadratinį arba kubinį keitimą.
\end{enumerate}

\textbf{Judėjimo uždavinys.}
\begin{enumerate}[label=\arabic*.]
\item Nubraižykite schemą: kryptys, keliai, laikai, greičiai.
\item Suvienodinkite vienetus.
\item Pasirinkite formulę \(s=vt\), \(v=\frac{s}{t}\) arba \(t=\frac{s}{v}\).
\item Dviejų kūnų uždaviniuose apsvarstykite bendrą greitį.
\end{enumerate}

\paragraph{Dažnos klaidos.}
Negalima maišyti minučių ir valandų toje pačioje formulėje. Jei greitis km/h, laikas turi būti valandomis. Keičiant plotą, \(1\text{ m}=100\text{ cm}\), bet \(1\text{ m}^2=100\cdot 100=10\,000\text{ cm}^2\), ne \(100\text{ cm}^2\).

\subsection{Iliustruojantys pavyzdžiai}
\textbf{1 pavyzdys.} Paverskime \(3,5\) km metrais.
\[
3,5\text{ km}=3,5\cdot 1000=3500\text{ m}.
\]
\textbf{Atsakymas:} \(3500\) m.

\textbf{2 pavyzdys.} Automobilis važiavo \(2,5\) h \(80\) km/h greičiu. Kokį kelią nuvažiavo?
\[
s=vt=80\cdot 2,5=200\text{ km}.
\]
\textbf{Atsakymas:} \(200\) km.

\textbf{3 pavyzdys.} Du dviratininkai išvažiavo vienas prieš kitą iš miestelių, tarp kurių \(54\) km. Jų greičiai \(12\) km/h ir \(15\) km/h. Po kiek laiko susitiks?

Artėjimo greitis:
\[
12+15=27\text{ km/h}.
\]
Laikas:
\[
t=\frac{54}{27}=2\text{ h}.
\]
\textbf{Atsakymas:} po \(2\) h.

\paragraph{Pasitikrinkite.}
\begin{enumerate}[label=\arabic*.]
\item Paverskite \(450\) cm metrais.
\item Paverskite \(2\text{ m}^2\) kvadratiniais centimetrais.
\item Kiek laiko truks \(18\) km kelionė \(6\) km/h greičiu?
\end{enumerate}

\subsection{Ryšys su kitomis temomis}
Matavimo vienetai reikalingi geometrijai \ref{sec:plotas-pavirsius-ir-turis}, finansų tarifams, proporcijoms ir tiesiniams sąryšiams. Kelio formulė yra praktinis tiesioginio proporcingumo pavyzdys.

\section{Plokštumos ir erdvės figūrų konstravimas}\label{sec:plokstumos-ir-erdves-figuru-konstravimas}

\subsection{Apibrėžtys}
\begin{apibreztis}
Konstravimas yra figūros braižymas pagal nurodytas sąlygas naudojant liniuotę, skriestuvą, matlankį arba skaitmenines priemones.
\end{apibreztis}
\begin{apibreztis}
Lygios figūros sutampa uždėtos viena ant kitos. Panašios figūros turi lygius atitinkamus kampus, o atitinkamų kraštinių santykiai yra vienodi.
\end{apibreztis}
\begin{apibreztis}
Mastelis yra brėžinio ilgio ir tikrojo ilgio santykis. Pavyzdžiui, mastelis \(1:100\) reiškia, kad \(1\) cm brėžinyje atitinka \(100\) cm tikrovėje.
\end{apibreztis}
\begin{apibreztis}
Išklotinė yra plokščias erdvės kūno sienų išdėstymas, iš kurio galima sulankstyti kūną.
\end{apibreztis}

\subsection{Teoremos ir savybės}
\begin{savybe}
Trikampis egzistuoja, jei bet kurių dviejų jo kraštinių suma didesnė už trečiąją kraštinę.
\end{savybe}
\begin{savybe}
Jei dviejų trikampių trys atitinkamos kraštinės lygios, trikampiai yra lygūs.
\end{savybe}
\begin{savybe}
Panašiųjų trikampių atitinkamų kraštinių santykiai lygūs tam pačiam skaičiui, vadinamam panašumo koeficientu.
\end{savybe}
\begin{savybe}
Postūmis, posūkis ir atspindys išlaiko figūros dydį ir formą, o didinimas arba mažinimas pakeičia dydį, bet panašumo atveju išlaiko formą.
\end{savybe}

\subsection{Taisyklės ir algoritmai}
\textbf{Trikampio konstravimas pagal tris kraštines.}
\begin{enumerate}[label=\arabic*.]
\item Nubraižykite vieną kraštinę.
\item Iš vieno jos galo skriestuvu pažymėkite lanką, kurio spindulys lygus antrai kraštinei.
\item Iš kito galo pažymėkite lanką, kurio spindulys lygus trečiai kraštinei.
\item Lankų susikirtimo tašką sujunkite su pradinės kraštinės galais.
\end{enumerate}

\textbf{Mastelio uždavinys.}
\begin{enumerate}[label=\arabic*.]
\item Perskaitykite mastelį kaip santykį.
\item Suvienodinkite vienetus.
\item Sudarykite proporciją tarp brėžinio ir tikrovės ilgių.
\item Raskite nežinomą ilgį.
\end{enumerate}

\paragraph{Dažnos klaidos.}
Ne kiekvienos trys atkarpos sudaro trikampį: \(2\), \(3\), \(6\) cm netinka, nes \(2+3<6\). Mastelyje reikia atidžiai skirti brėžinio ir tikrovės ilgius. Išklotinėje sienos turi jungtis taip, kad sulanksčius jos nesidubliuotų ir neliktų spragų.

\subsection{Iliustruojantys pavyzdžiai}
\textbf{1 pavyzdys.} Ar galima nubraižyti trikampį, kurio kraštinės \(4\) cm, \(5\) cm ir \(10\) cm?

Tikriname:
\[
4+5=9<10.
\]
Dviejų trumpesnių kraštinių suma mažesnė už trečiąją, todėl trikampis neegzistuoja.

\textbf{Atsakymas:} negalima.

\textbf{2 pavyzdys.} Plane mastelis \(1:500\). Takelis plane yra \(7\) cm. Koks tikrasis ilgis?
\[
7\cdot 500=3500\text{ cm}=35\text{ m}.
\]
\textbf{Atsakymas:} \(35\) m.

\textbf{3 pavyzdys.} Panašių trikampių atitinkamos kraštinės yra \(6\) cm ir \(9\) cm. Mažesnio trikampio kita kraštinė \(8\) cm. Raskime atitinkamą didesnio trikampio kraštinę.

Panašumo koeficientas:
\[
k=\frac{9}{6}=1,5.
\]
Todėl
\[
8\cdot 1,5=12\text{ cm}.
\]
\textbf{Atsakymas:} \(12\) cm.

\paragraph{Pasitikrinkite.}
\begin{enumerate}[label=\arabic*.]
\item Ar egzistuoja trikampis su kraštinėmis \(5\), \(7\), \(9\) cm?
\item Mastelis \(1:200\), brėžinyje ilgis \(4\) cm. Koks tikrasis ilgis metrais?
\item Panašumo koeficientas \(3\). Mažesnė kraštinė \(5\) cm. Kokia atitinkama didesnė kraštinė?
\end{enumerate}

\subsection{Ryšys su kitomis temomis}
Konstravimas jungia kampus \ref{sec:kampai-trikampiai-ir-keturkampiai}, proporcijas, mastelį ir erdvės kūnų paviršiaus skaičiavimą. Braižymas taip pat lavina matematinį argumentavimą: reikia paaiškinti, kodėl gauta figūra atitinka sąlygą.

\section{Duomenų tyrimas ir skaitmeninės technologijos}\label{sec:duomenu-tyrimas-ir-skaitmenines-technologijos}

\subsection{Apibrėžtys}
\begin{apibreztis}
Statistinis klausimas yra klausimas, į kurį atsakant tikimasi duomenų įvairovės, pavyzdžiui, „Kiek laiko klasės mokiniai kasdien skaito?“
\end{apibreztis}
\begin{apibreztis}
Kokybiniai duomenys apibūdina rūšį arba savybę, pavyzdžiui, spalvą. Kiekybiniai duomenys yra skaitinės reikšmės, pavyzdžiui, laikas minutėmis.
\end{apibreztis}
\begin{apibreztis}
Skaitmeninė priemonė yra programa arba įrankis, padedantis rinkti, tvarkyti, skaičiuoti ir vaizduoti duomenis.
\end{apibreztis}
\begin{apibreztis}
Duomenų interpretavimas yra duomenimis pagrįstos išvados formulavimas, nurodant, ką duomenys rodo ir ko jie nerodo.
\end{apibreztis}

\subsection{Teoremos ir savybės}
\begin{savybe}
Geras statistinis klausimas turi aiškų tiriamą požymį, aiškią tiriamą grupę ir numato galimus skirtingus atsakymus.
\end{savybe}
\begin{savybe}
Skaitmeninės priemonės greitai skaičiuoja ir braižo diagramas, bet įvestų duomenų teisingumą turi patikrinti žmogus.
\end{savybe}
\begin{savybe}
Dviguba stulpelinė diagrama tinka dviem grupėms palyginti pagal tas pačias kategorijas, o linijinė diagrama tinka reikšmės kitimui laike vaizduoti.
\end{savybe}
\begin{savybe}
Išvada turi atitikti imtį: jei apklausta tik viena klasė, negalima užtikrintai teigti apie visos šalies mokinius.
\end{savybe}

\subsection{Taisyklės ir algoritmai}
\textbf{Duomenų tyrimo planas.}
\begin{enumerate}[label=\arabic*.]
\item Suformuluokite statistinį klausimą.
\item Nuspręskite, kokius duomenis rinksite: kokybinius ar kiekybinius.
\item Parenkite aiškius apklausos arba stebėjimo klausimus.
\item Surinkite duomenis ir patikrinkite, ar nėra akivaizdžių klaidų.
\item Sudarykite lentelę, diagramą ir apskaičiuokite tinkamas charakteristikas.
\item Parašykite išvadą, remdamiesi skaičiais.
\end{enumerate}

\textbf{Darbas skaičiuoklėje.}
\begin{enumerate}[label=\arabic*.]
\item Kiekvieną požymį rašykite atskirame stulpelyje.
\item Vienoje eilutėje laikykite vieno stebinio duomenis.
\item Vienetus nurodykite stulpelio pavadinime, o langeliuose rašykite tik skaičius.
\item Prieš braižydami diagramą patikrinkite, ar pažymėjote tinkamą duomenų sritį.
\end{enumerate}

\paragraph{Dažnos klaidos.}
Apklausos klausimas neturi sufleruoti atsakymo. Diagrama be mastelio arba su nukirsta ašimi gali klaidinti. Skaitmeninė priemonė gali teisingai apskaičiuoti neteisingai įvestus duomenis, todėl pirmiausia reikia patikrinti duomenų prasmę.

\subsection{Iliustruojantys pavyzdžiai}
\textbf{1 pavyzdys.} Ar klausimas „Ar tau patinka mūsų nuostabi mokyklos valgykla?“ yra geras statistinis klausimas?

Klausimas sufleruoja teigiamą nuomonę žodžiu „nuostabi“. Geriau klausti: „Kaip įvertintum mokyklos valgyklą nuo \(1\) iki \(5\)?“

\textbf{Atsakymas:} pradinis klausimas netinkamas, nes yra šališkas.

\textbf{2 pavyzdys.} Kuri diagrama tiktų parodyti, kaip savaitės dienomis keitėsi klasės surinktų žingsnių skaičius?

Kadangi duomenys kinta laike, tinka linijinė diagrama. Ji parodo didėjimą, mažėjimą ir pokyčių kryptį.

\textbf{Atsakymas:} linijinė diagrama.

\textbf{3 pavyzdys.} Apklausėme \(20\) mokinių, kiek minučių jie sportuoja per dieną. Duomenų vidurkis \(38\) min, mediana \(30\) min. Kaip interpretuoti?

Vidurkis didesnis už medianą, vadinasi, tikėtina, kad keli mokiniai sportuoja gerokai ilgiau ir pakelia vidurkį. Išvada turi būti apie apklaustus \(20\) mokinių.

\textbf{Atsakymas:} daugumos sportavimo laikas gali būti arčiau \(30\) min, bet keli dideli laikai padidino vidurkį.

\paragraph{Pasitikrinkite.}
\begin{enumerate}[label=\arabic*.]
\item Patobulinkite klausimą: „Ar sutinki, kad matematika lengva?“
\item Kokią diagramą rinktumėtės dviejų klasių mėgstamiausiems vaisiams palyginti?
\item Kodėl prieš skaičiuojant vidurkį reikia patikrinti, ar visi duomenys turi tuos pačius vienetus?
\end{enumerate}

\subsection{Ryšys su kitomis temomis}
Duomenų tyrimas apibendrina statistiką \ref{sec:statistiniai-duomenys-ir-vidurkiai}, diagramas, procentus ir skaitmeninį raštingumą. Jis moko atsakingai pagrįsti teiginius skaičiais ir pasiruošia sudėtingesniam statistiniam patikimumui vyresnėse klasėse.
